% ============================================================================
%  v1.tex -- Paper draft skeleton (structure only)
% ============================================================================

\documentclass[11pt]{article}

% ---- encoding & fonts ------------------------------------------------------
\usepackage[T1]{fontenc}
\usepackage{lmodern}

% ---- page geometry (A4, ~2.5 cm margins; adjust to your needs) ----
\usepackage[a4paper, margin=2.5cm]{geometry}

% ---- math ------------------------------------------------------------------
\usepackage{amsmath}
\usepackage{amssymb}
\usepackage{mathtools}

% ---- figures & tables ------------------------------------------------------
\usepackage{graphicx}
\usepackage{booktabs}
\usepackage{multirow}
\usepackage{array}
\usepackage{caption}
\usepackage{subcaption}   % subfigures, e.g., Figure 3A/B/C
\usepackage{float}
\usepackage{cite}
\graphicspath{{figures/}} % put your figures under figures/

% ---- misc ------------------------------------------------------------------
\usepackage{enumitem}
\usepackage{xcolor}
\usepackage[hidelinks]{hyperref}  % load last
\usepackage{indentfirst}    % indent the first paragraph after a heading
\setlength{\parindent}{2em}    % two-character paragraph indent
\emergencystretch=1.5em    % avoid overfull hbox in paragraphs


% ---- commonly used macros (keep or remove as needed) -----------------------
\newcommand{\eg}{e.g.,}
\newcommand{\ie}{i.e.,}
\newcommand{\etal}{et al.}
% Uncomment the line below if you want visible red TODO markers in the PDF:
% \newcommand{\todo}[1]{\textcolor{red}{\textbf{[TODO] #1}}}

% ---- make the Abstract title the same size as section headings --------------
% (\Large\bfseries, same as \section; abstract body text stays \small)
\makeatletter
\renewenvironment{abstract}{%
  \small
  \begin{center}%
    {\Large\bfseries\abstractname\vspace{-.5em}\vspace{\z@}}%
  \end{center}%
  \quotation
}{%
  \endquotation
}
\makeatother

% ============================================================================
%  Title / authors (TODO: fill in yourself)
% ============================================================================
\title{GHXTox: Peptide toxicity prediction by fusing a task-adaptive protein language model with pLDDT-aware multi-scale 3D structural modeling}
\author{LaLaLa, LaLaLa, LaLaLa, LaLaLa}
\date{}% 不显示日期；如需显示可改为 \date{\today}

\begin{document}

\maketitle

% ----------------------------------------------------------------------------
%  Abstract & keywords 
% ----------------------------------------------------------------------------
\begin{abstract}
As biomedical research advances from the cellular to the molecular level,
peptides have become increasingly important as a versatile class of bioactive
molecules and are now widely used in the targeted treatment of diseases such as
Parkinson's disease and cancer. However, toxicity assessment remains a
critical bottleneck in peptide-based drug development. Existing methods have
begun to explore multi-dimensional (multimodal) feature fusion and toxicity
prediction based on the three-dimensional (3D) conformations of peptides,
yet most of them still underutilize 3D structural information, and
systematic multi-dimensional fusion remains rare. Here, we propose GHXTox, a hierarchical dual-level fusion framework. In the
first level, a sequence-enhanced 3D structural expert integrates sequence
semantics with spatial structural features through predicted
local-distance difference test (pLDDT)-aware
cross-attention. In the second level, this structural expert is fused with a
task-adaptive one-dimensional (1D) sequence expert at the probability level,
using weights frozen on training out-of-fold (OOF) predictions. The 1D branch analyzes sequences
with a task-adaptive protein language model (ProtT5 fine-tuned with low-rank
adaptation, LoRA) and a bidirectional Mamba state-space model. The 3D expert
integrates an ESMFold-derived structural anchor, pLDDT-aware residue
geometry, full-atom-derived chemical
sites, and a source-aware, availability-masked conformer-statistics residual for peptides
with valid ensemble information; samples without conformer statistics fall
back exactly to the single-structure pathway. On the frozen historical
benchmark test sets, GHXTox achieves MCC values of 0.8726 and 0.6746 and
AUPRC values of 0.9675 and 0.8043 on Test1 and Test2, respectively. Frozen-model
analyses indicate that the conformer contribution is small and non-uniform and
that, under a frozen structural expert, substantial geometric variation does
not necessarily produce stable outcome-aligned probability variation. Further
analyses characterize residue-level sequence sensitivity,
confidence--gate associations, chemical-site perturbations, and expert
disagreement; these findings are interpreted as model-level evidence rather
than as causal biological mechanisms.
\end{abstract}

\noindent\textbf{Keywords:} peptide toxicity prediction; multimodal deep learning; protein language model;
bidirectional Mamba; pLDDT-aware multi-scale 3D modeling; interpretability
% ============================================================================
%  1 INTRODUCTION
% ============================================================================
\section{INTRODUCTION}

As biomedical research advances from the cellular to the molecular level,
peptides have become increasingly important as a versatile class of bioactive
molecules and are now widely used in the targeted treatment of diseases such as
Parkinson's disease and cancer \cite{muttenthaler2021}. Peptide toxicity is a critical safety
consideration throughout the drug development pipeline, and the accurate and
efficient identification of toxic peptides remains a major bottleneck in
peptide-based drug discovery \cite{zhang2025toxmsrc}. Because conventional experimental toxicity
assays are costly, time-consuming, and difficult to scale, reliable
computational predictors are highly desirable.

Computational peptide toxicity prediction has evolved from traditional machine
learning methods based on handcrafted physicochemical features toward deep
learning approaches with stronger representation-learning capability. Current
mainstream methods can be roughly grouped into two categories. The first
category operates on sequence-level (one-dimensional, 1D) and/or atomic-level
(two-dimensional, 2D) representations. For example, ToxMSRC combines
continuous bag-of-words (CBOW) embeddings, multi-scale convolutional neural
networks, and a bidirectional long short-term memory network with residual
connections, achieving strong performance on the benchmark dataset \cite{zhang2025toxmsrc}.
ToxPLTC integrates pooled ProtT5 embeddings, borderline-SMOTE augmentation,
and a text convolutional neural network \cite{liang2026toxpltc}. ToxiPep fuses context-aware
sequence representations with atomic-level graphs through a cross-attention
mechanism \cite{guan2025toxipep}. HyPepTox-Fuse further combines embeddings from multiple protein
language models with conventional descriptors \cite{tran2025hypeptoxfuse}. Although these models
achieve strong performance within their respective modalities, they largely
neglect the three-dimensional (3D) conformations of peptides, which are
closely related to toxicity; cross-dimensional information interaction and
mutual validation remain limited, and their interpretability is mostly
confined to the sequence level.

The second category attempts to exploit 3D structural information. PeptiTox
uses ESM2 sequence embeddings together with ESMFold-predicted 3D structures
and geometric graph neural networks \cite{wang2025peptitox}. StrucToxNet feeds ProtT5 sequence
embeddings and ESMFold-predicted structures into an equivariant graph neural
network (EGNN) and systematically analyzes structural quality \cite{jiao2025structoxnet}. These
studies demonstrate the value of 3D conformations for toxicity prediction.
However, such methods are constrained by the limited accuracy of current
structure prediction models on short peptides and typically rely on a single
predicted conformation; structural confidence, i.e., the predicted local
distance difference test (pLDDT) score, is either discarded through hard
filtering or ignored, and multi-conformer ensembles are not considered.
Moreover, the fusion between the sequence and structure branches remains
relatively shallow, and no systematic multi-dimensional fusion that combines a
task-adaptive sequence model with confidence-aware multi-scale structure
modeling has been established.

To address these limitations, we propose GHXTox, an interpretable hierarchical
dual-level fusion framework for peptide toxicity prediction. In the first
level, a sequence-enhanced 3D structural expert integrates sequence semantics
with spatial structural features through pLDDT-aware cross-attention. In the
second level, this structural expert is fused with a task-adaptive 1D sequence
expert at the probability level, with weights frozen on training out-of-fold
(OOF) predictions. In the 1D branch, a
task-adaptive protein language model (ProtT5 fine-tuned with low-rank
adaptation, LoRA) is combined with a bidirectional Mamba state-space model for
sequence analysis \cite{elnaggar2022prottrans,hu2022lora,gu2024mamba}. In the
3D branch, ESMFold provides a common single-structure anchor for all peptides
\cite{lin2023esmfold}. AlphaFold2/ColabFold, PepFlow, and BioEmu were examined
as conformer-generation references during structural development
\cite{jumper2021alphafold2,mirdita2022colabfold,abdin2024pepflow,lewis2025bioemu}.
The retained conformer-residual module was trained on 2{,}795 peptides with
validated ensemble statistics, comprising 761 PepFlow records of 11--15
residues and 2{,}034 AlphaFold2/ColabFold records of 16--30 residues. A
three-dimensional source indicator conditions the residual head, while a hard
availability mask preserves the single-structure output for all uncovered
records. Frozen benchmark inference used only audited PepFlow and
AlphaFold2/ColabFold conformer caches under source- and sequence-matching
rules fixed independently of the benchmark labels. BioEmu was not included in
the retained benchmark inputs. Gated sequence--structure cross-attention in the
adjacent neuropeptide-prediction task further informs the interaction design,
without transferring its labels or task-specific decision rule
\cite{zhang2026neuroppredpscg}.
On this basis, we design a multi-scale structural information extraction
strategy that progressively captures structural cues from local residue
functional groups (chemical sites) to the whole peptide, and the pLDDT
confidence output by structure prediction models is exploited through
continuous down-weighting and dynamic gating of structural features. The 1D ensemble
probability is first calibrated using training out-of-fold predictions, after
which the calibrated 1D probability and the structural-expert probability are
combined using OOF-frozen fusion weights. Ablation experiments show that the
contribution of structural evidence is dataset- and metric-dependent: fusion
improves several fixed-threshold and probability-quality measures on Test1
while maintaining broadly comparable performance on Test2. Post hoc analyses
characterize residue-level sensitivity, confidence gating, chemical-site
perturbations, and expert disagreement without assigning causal biological
meaning to these signals.

In summary, the main contributions of this work are threefold: (1) we propose
GHXTox, an interpretable hierarchical dual-level fusion framework that integrates a
task-adaptive 1D protein language model with pLDDT-aware multi-scale 3D
structure modeling; (2) the 3D branch designs a multi-scale structural
information extraction strategy spanning from local residue functional groups
to the whole peptide, and incorporates pLDDT confidence into structural
feature processing through continuous down-weighting and dynamic gating; and
(3) extensive experiments on the frozen historical benchmark test sets show
that GHXTox improves several fixed-threshold and probability-quality measures
on Test1 while maintaining broadly comparable performance on Test2, while
ablation, robustness, and interpretability analyses characterize the
conditions under which sequence and structural evidence are complementary.

% ============================================================================
%  2 MATERIALS AND METHODS
% ============================================================================
\section{MATERIALS AND METHODS}
% TODO: (optional) general remarks on the materials and methods.

\subsection{Datasets}
We evaluate GHXTox on the benchmark datasets constructed by CAPTP \cite{jiao2024captp} and
subsequently adopted by ToxMSRC \cite{zhang2025toxmsrc}. Following the data construction protocol
of ToxMSRC \cite{zhang2025toxmsrc}, toxic peptide sequences were merged from three databases,
namely CSM-Toxin \cite{morozov2023csmboxin}, ToxinPred2 \cite{sharma2022toxinpred2}, and ATSE \cite{wei2021atse}, and sequences longer
than 50 residues were removed. Non-toxic sequences were retrieved from UniProt
\cite{uniprot2019} using the query ``NOT KW-0800 AND NOT KW-0020 AND reviewed:true'', and
sequences containing non-standard amino acids or exceeding 50 residues were
excluded. After combining the two groups and removing samples with conflicting
labels, CD-HIT \cite{fu2012cdhit} was applied with a sequence-identity threshold of 0.9 to
eliminate redundant sequences, yielding 2{,}138 toxic and 5{,}375 non-toxic
peptides. Following the same partition as ToxMSRC \cite{zhang2025toxmsrc}, we used 85\% of the data
as the training set and the remaining 15\% as the benchmark Test1 set,
resulting in 1{,}818 toxic and
4{,}569 non-toxic training peptides and 320 toxic and 806 non-toxic Test1
peptides. An additional benchmark Test2 set, containing 46 toxic and 536
non-toxic peptides, is used for frozen historical evaluation. Both Test1 and
Test2 had been observed during historical project development and are
therefore treated as frozen historical benchmarks rather than newly blinded
external validation sets; their labels were not used in the formal OOF
calibration, fusion-weight selection, or threshold-freezing procedure.
ToxinPred~3.0 \cite{rathore2024toxinpred3} was consulted as an external method
and data-source reference during project development, but its records were not
added to the present training set or to the formal Test1/Test2 result tables.
Table~\ref{tab:datasets} summarizes the statistics of the datasets used in
this study.

% ---- Table 1: dataset statistics ------------------------------------------
\begin{table}[htbp]
    \centering
    \caption{Statistics of the datasets used in this study.}
    \label{tab:datasets}
    \begin{tabular}{lcccc}
        \toprule
        Dataset & Toxic & Non-toxic & Total & Use \\
        \midrule
        Training (benchmark) & 1,818 & 4,569 & 6,387 & Model training \\
        Test1 (benchmark)    &   320 &   806 & 1,126 & Benchmark test \\
        Test2                &    46 &   536 &   582 & Benchmark test \\
        \bottomrule
    \end{tabular}
\end{table}

\subsection{Overview of the Model}
Figure~\ref{fig:overview} shows an overview of GHXTox, a hierarchical
dual-level fusion framework for peptide toxicity prediction. The framework
consists of a task-adaptive one-dimensional (1D) sequence expert, a
sequence-enhanced three-dimensional (3D) structural expert, and a probability
fusion module driven by training out-of-fold (OOF) predictions. Within the
structural expert, GHXTox first performs feature-level interaction between
sequence semantics and spatial structural information through pLDDT-aware
cross-attention. At the top level of the model, the prediction probabilities
produced by the 1D sequence expert and the 3D structural expert are then
fused at the probability level. GHXTox therefore combines intra-expert
feature-level fusion with inter-expert probability-level fusion, preserving the
independent discriminative capability of each branch while exploiting
complementary information from amino acid sequences, spatial geometry, local
chemical environments, and conformational variability.

% ---- Figure 1: overall framework (placeholder) -----------------------------
\begin{figure}[htbp]
    \centering
    % TODO: \includegraphics[width=\linewidth]{figures/fig1_overview.pdf}
    \caption{Overall architecture of GHXTox. GHXTox adopts a hierarchical dual-level fusion
             design. The task-adaptive 1D sequence expert extracts toxicity-related
             sequence features with a LoRA-fine-tuned ProtT5 encoder and a bidirectional
             Mamba network, followed by three-seed ensembling and out-of-fold (OOF) Platt
             calibration to produce $p_{\mathrm{1D}}$. The sequence-enhanced 3D structural
             expert extracts multi-scale structural information through an ESM2-enhanced
             sequence encoder, a pLDDT-aware equivariant graph neural network, full-atom
             chemical-site modeling, and an availability-masked conformer residual branch.
             Within this expert, sequence and spatial residue representations are
             integrated through pLDDT-aware cross-attention to produce
             $p_{\mathrm{SE3D}}$. The two expert probabilities are fused with frozen
             weights of 0.72 and 0.28, and the final toxicity decision is obtained with a
             threshold of 0.4702225653.}
    \label{fig:overview}
\end{figure}

\subsubsection{Task-Adaptive 1D Sequence Expert}
The task-adaptive 1D sequence expert takes the amino acid sequence of a
peptide as input. After tokenization, the sequence is processed by the ProtT5
protein language model \cite{elnaggar2022prottrans} to obtain
contextualized, residue-level representations. To adapt this general-purpose
protein language model to peptide toxicity prediction, we employ low-rank
adaptation (LoRA) \cite{hu2022lora} for parameter-efficient
fine-tuning, learning task-specific updates while the majority of the
pretrained backbone parameters remain frozen. The resulting residue-level
representations are then fed into a bidirectional Mamba state-space network
\cite{gu2024mamba}, which models residue dependencies in both the
N-to-C and C-to-N directions. After residual connections, layer
normalization, and masked pooling, a classification head produces the
sequence-level toxicity probability. To reduce the variance arising from
random initialization, we train three models on the full training set with
different random seeds and average their output probabilities. The ensemble
probability is then calibrated with Platt scaling \cite{platt1999} fitted
exclusively on training out-of-fold (OOF) predictions, yielding the final 1D
expert probability $p_{\mathrm{1D}}$.

\subsubsection{Sequence-Enhanced 3D Structural Expert}
The sequence-enhanced 3D structural expert comprises four subbranches: an
internal sequence-semantic subbranch, a spatial-geometric subbranch, a
full-atom chemical-site subbranch, and an availability-masked conformer-ensemble
residual subbranch. Collectively, these components cover five complementary
representation scales: residue-level semantics, atom-level chemical sites,
spatial residue neighborhoods, peptide-level global properties, and
conformer-ensemble statistics. In the sequence-semantic subbranch, amino acid residue
types, functional-group properties, explicit physicochemical descriptors, and
ESM2 residue embeddings \cite{lin2023esmfold} are projected into a
shared feature space and then processed by a two-layer Transformer encoder
\cite{vaswani2017transformer} to generate contextualized sequence
representations. In parallel, the spatial-geometric subbranch constructs a
hybrid residue graph that combines sequential adjacency edges with
spatial-neighborhood edges derived from the predicted peptide structure.
Residue physicochemical properties, local geometric descriptors, ESM2
embeddings, three-dimensional coordinates, and residue-level pLDDT scores are
then fed into three pLDDT-aware equivariant graph neural network (EGNN)
\cite{satorras2021egnn} layers to update the residue states and
spatial representations.

To characterize local chemical environments that cannot be fully captured by
a C$\alpha$-based residue graph alone, we extract chemical sites corresponding to
charged, aromatic, and hydrophobic groups as well as hydrogen-bond donors and
acceptors from the full-atom peptide structure. Message passing
\cite{gilmer2017mpnn} is then performed over these sites using pairwise
distances, directional relationships, interaction types, and structural
confidence. The resulting chemical-site features are mapped back into the
corresponding residue-level representations through a zero-initialized
residual projection. This design allows the model to progressively learn
toxicity-related local spatial--chemical patterns while preserving the initial
geometric backbone. For peptides with available conformer ensembles, the
model further extracts residue-level root-mean-square fluctuation (RMSF),
pairwise distance variance, contact occupancy under multiple distance
cutoffs, and peptide-level root-mean-square deviation (RMSD) and radius of
gyration. These statistics are incorporated through an availability-masked
residual branch to refine the structural expert's output. When
conformer-ensemble information is unavailable, this residual contribution is
constrained to exactly zero.

Within the structural expert, the contextualized sequence and
confidence-aware spatial representations are integrated through pLDDT-aware
cross-attention, followed by peptide-level confidence gating; the detailed
construction is given in Section 2.4.6. After masked mean pooling, fusion with
global peptide-level physicochemical features, and multilayer perceptron
(MLP) classification, the sequence-enhanced structural expert obtains a base
structural logit. The base structural logit is further corrected by an
availability-masked, zero-initialized conformer residual before being
transformed into the structural-expert probability $p_{\mathrm{SE3D}}$.

\subsubsection{Hierarchical Fusion and Toxicity Decision}
At the top level of the model, the calibrated 1D sequence expert and the
sequence-enhanced 3D structural expert are not combined through further
intermediate-feature concatenation; instead, their outputs are fused with
fixed weights in the probability space. The final toxicity probability is
defined as
\begin{equation}
    p_{\mathrm{GHXTox}} = 0.72\,p_{\mathrm{1D}} + 0.28\,p_{\mathrm{SE3D}}.
    \label{eq:fusion}
\end{equation}
The larger sequence-expert weight is an empirical outcome of the OOF
fusion search rather than a manually assigned prior; the structural expert
provides complementary evidence on spatial geometry, local chemical
environments, and conformational dynamics. The fusion weights are determined
exclusively from group-aware training out-of-fold (OOF) predictions and are
frozen before evaluation on the frozen historical benchmark test sets. In the implementation, the two
expert probabilities are likewise combined linearly, without re-optimizing
the fusion parameters on the test data.

Finally, the fused probability is converted into a binary toxicity
prediction using a fixed decision threshold:
\begin{equation}
    \hat{y} = \begin{cases}
        1, & p_{\mathrm{GHXTox}} \geq 0.4702225653,\\
        0, & p_{\mathrm{GHXTox}} < 0.4702225653,
    \end{cases}
    \label{eq:decision}
\end{equation}
where $\hat{y}=1$ denotes a toxic peptide and $\hat{y}=0$ denotes a
non-toxic peptide. The probability-calibration parameters, expert-fusion
weights, and classification threshold are all determined solely from training
out-of-fold (OOF) predictions. The current procedure does not use Test1 or
Test2 labels to fit these quantities. Because both sets had nevertheless been
observed during historical project development, their results are reported as
frozen retrospective benchmarks rather than newly blinded external
validation. This hierarchical design enables fine-grained
sequence--structure interaction through pLDDT-aware cross-attention inside the
structural expert, while the frozen probability-level fusion at the top level
enhances inter-expert complementarity and predictive stability.

\subsection{Task-Adaptive 1D Sequence Modeling}
The 1D sequence expert is designed to learn toxicity-related contextual
semantics and sequential dependencies directly from peptide amino acid
sequences. Instead of reducing a frozen protein-language-model output to a
fixed mean vector, GHXTox preserves the residue-level representations
generated by ProtT5 and adapts the pretrained encoder to the peptide toxicity
prediction task through low-rank adaptation (LoRA). The resulting
residue-level representations are then processed by a bidirectional Mamba
state-space encoder to jointly model local residue patterns, N-to-C
dependencies, and C-to-N dependencies. The encoded residue features are
aggregated via masked mean pooling and passed to a linear classification head
to produce the single-model toxicity probability. Finally, three models
trained on the full training set with different random seeds are combined by
equal-weight probability averaging, followed by Platt calibration fitted
exclusively on training out-of-fold (OOF) predictions, yielding the final 1D
expert probability $p_{\mathrm{1D}}$.

This design organizes pretrained protein semantics, parameter-efficient task
adaptation, bidirectional sequence modeling, and probability calibration into
a coherent processing pipeline. While most of the pretrained knowledge in
ProtT5 is retained, LoRA enables task-specific adaptation of the residue
representations. The bidirectional Mamba encoder further reorganizes these
representations according to toxicity-related sequential patterns, whereas
multi-seed ensembling and out-of-fold calibration reduce training variance
and align the calibrated probability scale with that of the downstream
structural expert.

\subsubsection{Residue Tokenization and LoRA-Adapted ProtT5}
Given a peptide sequence of length $L$, we first convert its amino acid
residues into the space-separated format required by the ProtT5 tokenizer and
map them to token indices. Special tokens are inserted automatically, and
sequences of different lengths are padded to a uniform length within each
batch. The maximum token length is set to 128, which is sufficient for all
peptides used in this study. Under the current ProtT5 tokenizer configuration,
the first $L$ token positions correspond to the $L$ amino acid residues,
followed by the EOS token and, when required, padding tokens. Accordingly, the
implementation constructs the residue-validity mask directly from the
original peptide length:
\begin{equation}
\mathcal{S}=(a_1,a_2,\ldots,a_L), \qquad
m_i=
\begin{cases}
1, & 1\leq i\leq L,\\
0, & i>L.
\end{cases}
\label{eq:sequence_mask}
\end{equation}
where $m_i=1$ denotes an actual amino acid position and $m_i=0$ denotes a
special or padding position; the same mask is applied in the Mamba scans and
the subsequent global pooling.

The tokenized sequence is then processed by the ProtT5-XL-UniRef50 encoder
\cite{elnaggar2022prottrans} to generate contextualized residue-level
representations:
\begin{equation}
\mathbf H^{\mathrm{PLM}}
=
\operatorname{ProtT5}_{\mathrm{LoRA}}(\mathcal S),
\qquad
\mathbf H^{\mathrm{PLM}}
\in\mathbb R^{L\times1024}.
\label{eq:prott5_representation}
\end{equation}
Retaining the complete residue-level representation matrix enables the
downstream encoder to preserve position-specific contextual semantics. Each
residue vector thus encodes not only its amino acid identity but also its
conditional context within the complete peptide sequence.

To adapt the general-purpose ProtT5 encoder to the peptide toxicity prediction
task, we apply low-rank adaptation (LoRA) \cite{hu2022lora} to the
linear transformations within the encoder. For an original linear layer
$\mathbf y=\mathbf W_0\mathbf x+\mathbf b$, the pretrained weight matrix
$\mathbf W_0$ is frozen, and the task-specific update is factorized into two
trainable low-rank matrices:
\begin{equation}
\mathbf W
=
\mathbf W_0+
\frac{\alpha}{r}\mathbf B\mathbf A,
\qquad
\mathbf A\in\mathbb R^{r\times d_{\mathrm{in}}},
\quad
\mathbf B\in\mathbb R^{d_{\mathrm{out}}\times r}.
\label{eq:lora_update}
\end{equation}
The LoRA-augmented forward pass is therefore
\begin{equation}
\mathbf y
=
\mathbf W_0\mathbf x+\mathbf b
+
\frac{\alpha}{r}
\mathbf B\mathbf A
\operatorname{Dropout}(\mathbf x).
\label{eq:lora_forward}
\end{equation}
We set the LoRA rank to $r=8$, the scaling coefficient to $\alpha=16$, and
the LoRA dropout rate to $0.1$, yielding a scaling factor of $\alpha/r=2$.
Matrix $\mathbf A$ is initialized with Kaiming initialization, whereas
$\mathbf B$ is initialized to zero; consequently, the LoRA branch initially
leaves the pretrained transformation unchanged and progressively learns
toxicity-specific updates during training.

Unlike conventional configurations that apply LoRA only to selected attention
projections, the current implementation augments all linear layers within the
ProtT5 encoder. The pretrained base transformations and embeddings remain
frozen, whereas the low-rank parameters and the downstream sequence-
classification modules are trainable.

The 1024-dimensional ProtT5 representations are then projected into the
192-dimensional hidden space used by the Mamba encoder \cite{gu2024mamba}:
\begin{equation}
\mathbf X^{(0)}
=
\operatorname{Dropout}
\left(
\operatorname{GELU}
\left[
\operatorname{LN}
\left(
\mathbf H^{\mathrm{PLM}}\mathbf W_p+\mathbf b_p
\right)
\right]
\right),
\qquad
\mathbf X^{(0)}\in\mathbb R^{L\times192}.
\label{eq:prott5_projection}
\end{equation}
where $\operatorname{LN}$ denotes layer normalization. The projection block
consists of a linear transformation, layer normalization, GELU activation,
and dropout. It reduces the computational dimension for the downstream
state-space encoder while providing a task-specific nonlinear transformation
of the ProtT5 residue representations.

\subsubsection{Bidirectional Mamba Sequence Encoder}
Although ProtT5 provides contextualized residue representations, these
representations primarily encode general protein semantics learned during
large-scale pretraining. To further model the sequential dependencies
specifically associated with peptide toxicity, two bidirectional Mamba layers
are placed after the ProtT5 projection block. Mamba employs an input-dependent
selective state-space mechanism that dynamically controls which residue
information is retained, updated, or suppressed during sequential scanning
\cite{gu2024mamba}.

For a unidirectional Mamba scan, the input representation
$\mathbf X \in \mathbb{R}^{L \times d_{\mathrm{model}}}$ with
$d_{\mathrm{model}} = 192$ is first projected into a content branch and a
gating branch:
\begin{equation}
    [\mathbf U, \mathbf Z] = \mathbf X \mathbf W_{\mathrm{in}} + \mathbf b_{\mathrm{in}},
    \qquad \mathbf U, \mathbf Z \in \mathbb{R}^{L \times d_{\mathrm{inner}}},
    \quad d_{\mathrm{inner}} = 2 d_{\mathrm{model}} = 384,
    \label{eq:mamba_proj}
\end{equation}
i.e., the projection expands the feature dimension from 192 to 768 and splits
it into two branches, each of dimension $d_{\mathrm{inner}} = 384$. The
content branch is processed by a depthwise one-dimensional convolution with a
kernel width of 4, followed by a SiLU activation:
\begin{equation}
    \widetilde{\mathbf U} = \operatorname{SiLU}(\operatorname{DWConv1D}(\mathbf U)),
    \qquad \widetilde{\mathbf U}_t^{\mathrm{m}} = m_t \widetilde{\mathbf U}_t,
    \label{eq:mamba_conv}
\end{equation}
where $\widetilde{\mathbf U}_t^{\mathrm{m}}$ denotes the convolved feature
after applying the residue-validity mask. This local convolution extracts
short-range residue patterns before the state-space scan, enabling the module
to combine local sequence motifs with the longer-range context accumulated
through state recurrence. Input-dependent projections then generate the
discretization step size and the state-space parameters from the masked
feature:
\begin{equation}
    [\delta_t, \mathbf B_t, \mathbf C_t] = \operatorname{Linear}(\widetilde{\mathbf U}_t^{\mathrm{m}}),
    \qquad \Delta_t = \operatorname{Softplus}(\mathbf W_{\Delta}\delta_t + \mathbf b_{\Delta}).
    \label{eq:mamba_params}
\end{equation}
The state-transition matrix is parameterized to be negative and is
discretized according to the input-dependent step size:
\begin{equation}
    \mathbf A = -\exp(\mathbf A_{\log}),
    \qquad \bar{\mathbf A}_t = \exp(\Delta_t \mathbf A),
    \qquad \bar{\mathbf B}_t = \Delta_t \mathbf B_t,
    \label{eq:mamba_discretize}
\end{equation}
where $\mathbf A_{\log} \in \mathbb{R}^{d_{\mathrm{inner}} \times d_{\mathrm{state}}}$
is a learnable parameter matrix. The selective state recurrence is expressed
as
\begin{equation}
    \mathbf S_t = \bar{\mathbf A}_t \odot \mathbf S_{t-1} + \bar{\mathbf B}_t \odot \widetilde{\mathbf U}_t^{\mathrm{m}},
    \qquad Y_t = \sum_{n=1}^{d_{\mathrm{state}}} S_{t,n} C_{t,n},
    \label{eq:mamba_recurrence}
\end{equation}
where the state dimension is set to $d_{\mathrm{state}} = 16$. The state-space
output at position $t$ is subsequently combined with the gating branch and the
skip term:
\begin{equation}
    \mathbf O_t = \mathbf W_{\mathrm{out}}\left[\mathbf Y_t \odot
    \operatorname{SiLU}(\mathbf Z_t) + \mathbf D \odot \widetilde{\mathbf U}_t^{\mathrm{m}}\right].
    \label{eq:mamba_output}
\end{equation}
At invalid positions, $\widetilde{\mathbf U}_t^{\mathrm{m}}$ is zero and
$\Delta_t$ is also zeroed, preventing padding tokens from injecting
information into the recurrent states.

A unidirectional scan accesses only one side of the sequence context, whereas
peptide toxicity may depend on residue patterns associated with both termini
and with interactions spanning different sequence positions. Therefore, in the
$\ell$-th bidirectional Mamba layer, forward and backward scans are applied to
the normalized sequence using the residue-validity mask $\mathbf m$:
\begin{equation}
    \begin{aligned}
        \mathbf F^{(\ell)} &= \mathcal{M}_f^{(\ell)}(\operatorname{LN}(\mathbf X^{(\ell)}), \mathbf m),\\
        \mathbf R^{(\ell)} &= \operatorname{Reverse}\left[\mathcal{M}_b^{(\ell)}
        \left(\operatorname{Reverse}\left[\operatorname{LN}(\mathbf X^{(\ell)})\right],
        \operatorname{Reverse}(\mathbf m)\right)\right],
    \end{aligned}
    \label{eq:mamba_bidir}
\end{equation}
where the backward-scan output is restored to the original residue order. The
two directional outputs, each of dimension 192, are concatenated along the
feature dimension (384 in total) and fused by a linear layer that projects the
representation back to 192 dimensions:
\begin{equation}
    \mathbf M^{(\ell)} = \operatorname{GELU}\left(\operatorname{Dropout}
    \left(\mathbf W_m \left[\mathbf F^{(\ell)} \parallel \mathbf R^{(\ell)}\right]
    + \mathbf b_m\right)\right),
    \qquad \mathbf X^{(\ell+1)} = \mathbf X^{(\ell)} + \operatorname{Dropout}(\mathbf M^{(\ell)}),
    \label{eq:mamba_fuse}
\end{equation}
where $\parallel$ denotes feature concatenation, the dropout inside the GELU
block is applied within the fusion module, and the second dropout is applied
before the residual addition. Two bidirectional Mamba layers are stacked in
the final model. The first layer reconstructs bidirectional residue
dependencies from the task-adapted ProtT5 features, whereas the second layer
further aggregates the resulting contextual representations. This design
allows each residue position to receive conditional information from both its
N-terminal and C-terminal contexts.

\subsubsection{Masked Pooling and Class-Weighted Classification}
After two bidirectional Mamba layers, the model produces residue-level
sequence representations:
\begin{equation}
    \mathbf X^{(2)} = [\mathbf x_1, \mathbf x_2, \ldots, \mathbf x_L],
    \qquad \mathbf X^{(2)} \in \mathbb{R}^{L \times 192}.
    \label{eq:repr}
\end{equation}
Because peptide sequences have variable lengths, GHXTox does not rely on a
fixed-position classification token. Instead, masked mean pooling is applied
over the valid amino acid positions:
\begin{equation}
    \mathbf h_{1D} = \frac{\sum_{i=1}^{L} m_i \mathbf x_i}{\sum_{i=1}^{L} m_i},
    \label{eq:masked_pool}
\end{equation}
where the denominator equals the number of valid residue positions. This
operation converts the variable-length residue sequence into a fixed
192-dimensional peptide-level representation while ensuring that EOS and
padding positions do not contribute to the pooled feature. In contrast to
global max pooling, masked mean pooling aggregates distributed evidence across
the entire peptide and avoids a representation dominated by a single extreme
activation.

The peptide-level representation is first normalized and then passed through
a single-output linear classification head:
\begin{equation}
    z_{1D} = \mathbf w_c^{\top} \operatorname{LN}(\mathbf h_{1D}) + b_c,
    \qquad p_{1D}^{\mathrm{raw}} = \sigma(z_{1D})
    = \frac{1}{1 + \exp(-z_{1D})},
    \label{eq:logit}
\end{equation}
where $z_{1D}$ denotes the 1D expert logit and $p_{1D}^{\mathrm{raw}}$ the
corresponding uncalibrated toxicity probability. The classification head of
the current 1D expert consists of layer normalization followed by a single
linear layer, rather than a multilayer perceptron; it is therefore referred to
as a linear classification head throughout the architecture and Methods
sections.

Because non-toxic peptides substantially outnumber toxic peptides in the
training data, a positive-class-weighted binary cross-entropy loss is used in
the current 1D expert instead of SMOTE or Borderline-SMOTE. For a training
partition containing $N_+$ toxic and $N_-$ non-toxic peptides, the
positive-class weight is defined as
\begin{equation}
    w_+ = \frac{N_-}{N_+},
    \label{eq:pos_weight}
\end{equation}
and the class-weighted objective is expressed as
\begin{equation}
    \mathcal{L}_{\mathrm{WBCE}} = -\frac{1}{N}
    \sum_{j=1}^{N} \left[ w_+ y_j \log \sigma(z_j)
    + (1 - y_j) \log(1 - \sigma(z_j)) \right].
    \label{eq:wbce}
\end{equation}
During five-fold out-of-fold training, $w_+$ is calculated independently from
the class counts of each fold's training partition; for full-data training, it
is recalculated using the complete training set. This strategy increases the
contribution of the minority toxic class without generating synthetic peptide
sequences or feature vectors, while preventing any information from the
held-out validation fold from entering the training procedure.

\subsubsection{Three-Seed Ensemble and OOF Platt Calibration}
The predictions of deep neural networks may vary because of parameter
initialization, mini-batch ordering, and stochastic optimization. To reduce
the variance associated with an individual training run, three 1D expert
models are trained under identical data, architecture, and optimization
settings using the random seeds $42$, $123$, and $2025$. The deployed ensemble
consists of three models trained on the complete training set, rather than the
fold-specific checkpoints produced during cross-validation.

For sample $j$, the toxicity probability produced by the model trained with
seed $k$ is denoted by $p_j^{(k)} = \sigma(z_j^{(k)})$, $k \in \{1,2,3\}$, and
the uncalibrated ensemble score is obtained by equal-weight arithmetic
averaging:
\begin{equation}
    s_j = \frac{1}{3} \sum_{k=1}^{3} p_j^{(k)}.
    \label{eq:ensemble_score}
\end{equation}
No individual seed is selected according to its performance on Test1 or
Test2, and no seed-specific ensemble weight is introduced. All members share
the same architecture and training protocol; the different random seeds
provide averaging over a limited set of initialization and optimization
stochasticities.

Although probability averaging reduces prediction variance, neural-network
probabilities may remain systematically overconfident or underconfident.
Moreover, the top-level GHXTox framework subsequently combines the 1D expert
with the structural expert, making probability-scale compatibility
particularly important. Therefore, Platt scaling \cite{platt1999} is applied
after the three-model probability average.

Training out-of-fold predictions are first generated using the fixed
group-aware five-fold partition. Each training peptide receives an OOF
probability from a model trained without the fold containing that peptide.
The OOF probabilities from the three random seeds are then averaged to obtain
the OOF ensemble score $s_j^{\mathrm{OOF}}$. A two-parameter logistic
calibrator is fitted using only these training OOF scores and their labels:
\begin{equation}
    p_{1D,j} = \sigma\left(a\, \mathrm{logit}\left(\widetilde{s}_j\right) + b\right),
    \label{eq:platt_form}
\end{equation}
where $\mathrm{logit}(s) = \ln(s/(1-s))$,
$\widetilde{s}_j = \mathrm{clip}(s_j, \varepsilon, 1-\varepsilon)$ with
$\varepsilon = 10^{-6}$, and $a$ and $b$ denote the Platt calibration
parameters; the clip prevents infinite logits for probabilities of exactly
0 or 1, and the same clipping is applied to the OOF scores used in fitting. The parameters are estimated as
\begin{equation}
    \begin{aligned}
        (a, b) = \arg\min_{a,b} \;& \sum_{j=1}^{N}
        \mathrm{BCE}\left[y_j,\; \sigma\left(a\,
        \mathrm{logit}\left(\widetilde{s}_j^{\mathrm{OOF}}\right) + b\right)\right].
    \end{aligned}
    \label{eq:platt_fit}
\end{equation}
The scaling parameter is reparameterized as $a = \mathrm{clip}(\exp(\eta),\, 0.02, 50) > 0$, which
guarantees a monotonic probability transformation and therefore preserves the
sample ranking of the ensemble score; AUROC and AUPRC are consequently
unchanged by calibration. For benchmark test or previously unseen peptides,
the output probabilities of the three full-data models are first averaged,
after which the frozen calibration parameters are applied:
\begin{equation}
    p_{\mathrm{1D}} = \sigma\left(a\, \mathrm{logit}\left(\mathrm{clip}
    \left(\frac{1}{3} \sum_{k=1}^{3} p^{(k)},\; \varepsilon,\;
    1-\varepsilon\right)\right) + b\right).
    \label{eq:platt_infer}
\end{equation}
The labels of Test1 and Test2 are not used for ensemble-member selection,
calibration fitting, or 1D threshold optimization. The calibrated probability
$p_{\mathrm{1D}}$ is the sequence-expert output used in the top-level
probability fusion described in Section 2.2.3. The complete order of
operations is therefore
\begin{equation*}
    \text{3-seed probabilities} \rightarrow \text{averaging} \rightarrow
    \text{OOF-fitted Platt calibration} \rightarrow p_{\mathrm{1D}} \rightarrow
    \text{1D/3D fusion}.
\end{equation*}
Calibration is performed within the 1D expert before the top-level 1D/3D
fusion; no additional Platt calibrator is fitted after the final GHXTox
probability fusion.

\subsection{pLDDT-Aware Multi-Scale 3D Structural Modeling}
Peptide toxicity is influenced not only by amino acid composition and linear
residue order but also by the three-dimensional (3D) arrangement of charged
groups, hydrophobic regions, aromatic systems, and sulfur-containing
moieties. However, peptides are generally short, conformationally flexible,
and poorly covered by experimentally determined structures. Treating a single
predicted conformation as an error-free structural observation may therefore
amplify structure-prediction errors. Conversely, representing a peptide with
only C$\alpha$ coordinates or peptide-level summary descriptors fails to capture the
local spatial-chemical environments that are relevant to toxicity.

To address these limitations, GHXTox employs a pLDDT-aware multi-scale
structural expert that captures complementary information at five levels:
residue semantics, spatial residue neighborhoods, full-atom chemical sites,
peptide-level global properties, and conformer-ensemble statistics. The model
first extracts residue-level C$\alpha$ coordinates, full-atom coordinates, and
normalized pLDDT confidence scores from the predicted structure. An internal
sequence-semantic subbranch combines amino acid residue types,
functional-group categories, explicit physicochemical descriptors, and ESM2
residue embeddings to generate contextualized sequence states. In parallel, a
spatial-geometric subbranch constructs a hybrid residue graph containing both
sequential adjacency edges and spatial-neighborhood edges and propagates
geometric information through pLDDT-aware equivariant graph layers.

To overcome the limited chemical resolution of a C$\alpha$-only residue graph,
GHXTox additionally extracts typed chemical sites from the full-atom
structure and performs message passing using dual-scale distances, directional
relationships, interaction-compatibility channels, and structural confidence.
For peptides with available conformer ensembles, an additional statistical
residual branch summarizes residue fluctuations, distance variation, and
contact occupancy. The residual is constrained to exactly zero when conformer
information is unavailable.

Rather than discarding low-confidence structures through hard filtering,
GHXTox treats pLDDT as a continuous reliability signal and incorporates it
into geometric edge weighting, chemical-site messages, residue-level
structural gating, and peptide-level structural gating. Low-confidence
structural evidence is therefore down-weighted rather than removed, preserving
structural coverage while limiting the influence of uncertain coordinates.
Importantly, pLDDT is interpreted here as a local confidence estimate produced
by the structure-prediction model, rather than an experimentally determined
measure of structural accuracy, conformational flexibility, or biological
activity. The detailed architecture of the structural expert is illustrated in
Figure~\ref{fig:expert_detail}.

% ---- Figure 2: detailed 3D expert architecture (placeholder) ---------------
\begin{figure}[htbp]
    \centering
    % TODO: \includegraphics[width=\linewidth]{figures/fig2_expert_detail.pdf}
    \caption{Detailed architecture of the sequence-enhanced structural expert.
             The structural expert integrates ESM2-enhanced residue semantics, a
             pLDDT-aware geometric residue graph, full-atom chemical-site message
             passing, and availability-masked conformer-ensemble statistics.
             Residue-level structural states are continuously modulated by
             node-level confidence before serving as keys and values in
             sequence--structure cross-attention. A peptide-level confidence gate
             further controls the fused structural context, while conformer
             statistics provide a zero-initialized logit residual only when
             valid ensemble information is available.}
    \label{fig:expert_detail}
\end{figure}

\subsubsection{Peptide Structure Prediction and Conformer Preparation}
For a peptide sequence of length $L$, $S = (a_1, a_2, \ldots, a_L)$,
ESMFold is first used to predict its three-dimensional structure
\cite{lin2023esmfold}. The predicted structure is stored as a full-atom PDB file,
from which residue-level C$\alpha$ coordinates, backbone coordinates,
full-atom-derived chemical-site coordinates, and pLDDT scores are extracted.
This operation is denoted as
\begin{equation}
    (\mathbf X, \mathcal{A}, \mathbf q) = \operatorname{ESMFold}(\mathcal S),
    \qquad \mathbf X = [\mathbf x_1, \ldots, \mathbf x_L]^{\top} \in \mathbb{R}^{L \times 3},
    \quad \mathbf q = (q_1, \ldots, q_L) \in [0,1]^{L},
    \label{eq:esmfold}
\end{equation}
where $\mathbf X$ contains the C$\alpha$ coordinates, $\mathcal{A}$ denotes the
full-atom-PDB-derived backbone and chemical-site geometry, and $\mathbf q$ contains the
residue-level pLDDT confidence scores rescaled to $[0,1]$. The C$\alpha$ coordinates
serve as the primary spatial variables of the residue-level equivariant
graph, whereas the full-atom coordinates are used to construct local chemical
sites. Before geometric encoding, the coordinates are centered at the peptide
centroid to remove the effect of absolute translation.

A single predicted structure provides a common structural anchor but cannot
represent potential conformational heterogeneity. The retained
conformer-residual stage was therefore trained on a source-aware full2795
cohort containing 2{,}795 peptides with up to eight validated, aligned
conformers per peptide. PepFlow supplies 761 training records of 11--15
residues, whereas AlphaFold2/ColabFold supplies 2{,}034 records of 16--30
residues \cite{abdin2024pepflow,jumper2021alphafold2,mirdita2022colabfold}.
These sources are represented explicitly in the residual head; nevertheless,
generator identity remains partially coupled to peptide length and is treated
as a limitation rather than as an independently identified causal factor.

For frozen benchmark inference, conformer statistics were admitted only from
an audited cache whose source and sequence-matching policies were fixed without
using Test1 or Test2 labels. The Test1 cache contains eight conformers for each
of 486 peptides: 130 PepFlow-derived records of 3--15 residues and 356
AlphaFold2/ColabFold-derived records of 16--30 residues. Test2 reuses 284 of
these ensembles through an audited exact unique-sequence join; no independent
Test2 conformer generation or label-dependent cache selection is performed.
For every admitted ensemble, residue indices are matched to a common sequence
order. An ensemble is represented as
\begin{equation}
    \mathcal{C} = \{\mathbf X^{(1)}, \mathbf X^{(2)}, \ldots, \mathbf X^{(M)}\},
    \qquad 1 \leq M \leq 8, \quad \mathbf X^{(m)} \in \mathbb{R}^{L \times 3},
    \label{eq:conformer_set}
\end{equation}
where $\mathbf X^{(m)}$ denotes the C$\alpha$ coordinates of conformer $m$. Rigid-body
alignment is performed before the ensemble statistics are calculated,
ensuring that the resulting variation reflects internal geometric changes
rather than arbitrary global rotations or translations.

Because conformer ensembles are not available for every training peptide, a
binary availability indicator is defined as
\begin{equation}
    c_{\mathrm{conf}} =
    \begin{cases}
        1, & \text{if valid ensemble statistics are available,}\\
        0, & \text{otherwise.}
    \end{cases}
    \label{eq:conf_available}
\end{equation}
When $c_{\mathrm{conf}} = 0$, the conformer-statistics tensors are filled
with zeros, and the corresponding residual contribution is explicitly
suppressed by the availability indicator. Missing ensemble information is
therefore not interpreted as evidence of structural rigidity; instead, the
model falls back exactly to the single-structure pathway. Under the audited
benchmark cache, valid conformer statistics are available for 486/1{,}126
(43.16\%) Test1 peptides and 284/582 (48.80\%) Test2 peptides. Only these
covered records may receive a conformer residual; all remaining records retain
the exact single-structure prediction.

The generated conformers are computational samples rather than experimentally
determined trajectories or rigorous thermodynamic ensembles. Consequently,
RMSF, contact occupancy, and distance variance are interpreted as
sampling-derived descriptors of conformational variation rather than direct
measurements of peptide dynamics in solution.

\subsubsection{Residue Feature Construction and Hybrid Geometric Graph}
The structural expert constructs both an internal sequence representation and
a spatial-geometric representation at the residue level. Importantly, the
internal sequence-semantic subbranch described here uses frozen ESM2 residue
embeddings and a lightweight Transformer. It is distinct from the
LoRA-ProtT5 and bidirectional Mamba expert described in Section 2.3. The
former supports residue-level sequence--structure alignment inside the
structural expert, whereas the latter acts as an independent top-level 1D
prediction expert.

For residue $i$, the input to the internal sequence-semantic subbranch is
\begin{equation}
    \mathbf u_i^{\mathrm{seq}} = \left[\mathbf e_i^{\mathrm{AA}} \parallel
    \mathbf e_i^{\mathrm{group}} \parallel \mathbf f_i^{\mathrm{res}} \parallel
    P_{\mathrm{ESM}}\left(\mathbf e_i^{\mathrm{ESM2}}\right)\right],
    \label{eq:seq_feat}
\end{equation}
where $\mathbf e_i^{\mathrm{AA}} \in \mathbb{R}^{32}$ is the amino acid
category embedding, $\mathbf e_i^{\mathrm{group}} \in \mathbb{R}^{12}$ is the
functional-group category embedding, $\mathbf f_i^{\mathrm{res}} \in
\mathbb{R}^{42}$ denotes the explicit residue-level physicochemical
descriptors, and $\mathbf e_i^{\mathrm{ESM2}} \in \mathbb{R}^{1280}$ is the
ESM2-650M residue representation; the projection $P_{\mathrm{ESM}}$ maps the
ESM2 representation to 128 dimensions. Stacking the residue-level feature
vectors yields the sequence-feature matrix $U_{\mathrm{seq}} \in
\mathbb{R}^{L \times 214}$, where $214 = 32 + 12 + 42 + 128$. After a linear
projection and positional encoding, the sequence is processed by a two-layer
Transformer encoder with hidden dimension 128, four attention heads, and a
feed-forward dimension of 512:
\begin{equation}
    \mathbf H^{\mathrm{seq}} = \operatorname{Transformer}_2\left[
        P_{\mathrm{seq}}\left(U_{\mathrm{seq}}\right) + E_{\mathrm{pos}}
    \right], \qquad \mathbf H^{\mathrm{seq}} \in \mathbb{R}^{L \times 128},
    \label{eq:seq_transformer}
\end{equation}
This representation encodes the contextual sequence semantics of each residue
and subsequently serves as the query representation in the
sequence--structure cross-attention module.

In the spatial subbranch, the initial representation of residue $i$ is
constructed as
\begin{equation}
    \mathbf h_i^{(0)} = P_{\mathrm{sp}}\left[\mathbf f_i^{\mathrm{res}}
    \parallel \mathbf g_i^{\mathrm{geo}} \parallel q_i \parallel
    P_{\mathrm{ESM}}\left(\mathbf e_i^{\mathrm{ESM2}}\right)\right],
    \label{eq:node_init}
\end{equation}
where $\mathbf g_i^{\mathrm{geo}} \in \mathbb{R}^{16}$ is a lightweight
geometric descriptor derived from the C$\alpha$ coordinates and pLDDT scores. The
16-dimensional descriptor includes (i) residue-level pLDDT and local
confidence statistics; (ii) radial position, contact fractions, and soft
spatial density; (iii) nearest-neighbor, nonlocal-neighbor, and local
sequential distances; and (iv) sine/cosine encodings of the C$\alpha$-trace angles
and torsions. These explicit descriptors provide stable local geometric
priors without introducing an overly complex handcrafted representation for
short and uncertain predicted peptide structures.

GHXTox then constructs a hybrid residue graph:
\begin{equation}
    \mathcal{G} = \left(\mathcal{V}, \mathcal{E}_{\mathrm{seq}} \cup
    \mathcal{E}_{\mathrm{knn}}\right),
    \label{eq:hybrid_graph}
\end{equation}
where each residue corresponds to one node. Sequential edges
$\mathcal{E}_{\mathrm{seq}}$ connect adjacent residues and preserve
peptide-chain continuity, whereas $\mathcal{E}_{\mathrm{knn}}$ connects each
residue to its 10 nearest spatial neighbors and is symmetrized. The hybrid
graph therefore captures both covalent sequence order and nonlocal spatial
contacts while avoiding large variations in graph density that may arise from
a single fixed-distance cutoff.
 Equivalently, $A_{ij} \in \{0,1\}$ denotes whether residues $i$ and $j$ are
connected by a sequential or k-nearest-neighbor edge in the hybrid graph.

For residues $i$ and $j$, the C$\alpha$ distance is
\begin{equation}
    d_{ij} = \lVert \mathbf x_i - \mathbf x_j \rVert_2,
    \qquad \phi_k(d_{ij}) = \exp\left[-\gamma\left(d_{ij} - \mu_k\right)^2\right],
    \quad \gamma = \frac{1}{(\mu_{k+1} - \mu_k)^2},
    \quad k = 1, \ldots, 16,
    \label{eq:rbf}
\end{equation}
where the centers $\mu_k$ are uniformly spaced over the range 0--24 Å. The
base edge representation additionally includes the pairwise structural
confidence and the normalized sequence separation:
\begin{equation}
    c_{ij} = \sqrt{q_i q_j}, \qquad
    s_{ij} = \frac{\log(1 + |i - j|)}{\log(L + 1)}.
    \label{eq:edge_feat}
\end{equation}
Thus, sequentially adjacent residues receive small values, whereas residues
separated by longer sequence intervals approach one. The resulting edge
feature is therefore $\mathbf e_{ij} = \left[\boldsymbol\phi(d_{ij}) \parallel
c_{ij} \parallel s_{ij}\right]$. This formulation allows the model to
distinguish sequentially adjacent contacts, long-range spatial contacts, and
contacts supported by different levels of structural confidence.

\subsubsection{pLDDT-Aware Equivariant Graph Encoding}
To learn three-dimensional representations without relying on a particular
global coordinate frame, the spatial branch employs three E(3)-equivariant
graph neural network layers \cite{satorras2021egnn}. The network updates residue features
using rotation- and translation-invariant quantities while updating
coordinates equivariantly. Consequently, a global rigid transformation of the
peptide structure does not change the predicted toxicity, whereas
intermediate coordinate representations transform consistently with the
input.

At layer $\ell$, the message from residue $j$ to residue $i$ is computed as
\begin{equation}
    \mathbf m_{ij}^{(\ell)} = \phi_e^{(\ell)}\left[\mathbf h_i^{(\ell)}
    \parallel \mathbf h_j^{(\ell)} \parallel \mathbf e_{ij}^{(\ell)}\right],
    \label{eq:egnn_message}
\end{equation}
where $\phi_e^{(\ell)}$ is a message multilayer perceptron. To prevent
uncertain or distant contacts from contributing as strongly as local
high-confidence contacts, a continuous edge weight is defined as
\begin{equation}
    w_{ij}^{(\ell)} = A_{ij}\, \exp\!\left(-\frac{d_{ij}^{(\ell)}}{12}\right)
    \left[q_{\min} + \left(1 - q_{\min}\right) c_{ij}^{\rho}\right],
    \qquad q_{\min} = 0.1, \quad \rho = 1.5,
    \label{eq:edge_weight}
\end{equation}
The term $\exp(-d_{ij}^{(\ell)}/12)$ imposes a smooth distance decay with
a characteristic decay length of 12 Å, and both the distance and the edge
weight are recomputed at each layer from the current coordinates. Here,
$A_{ij} \in \{0,1\}$ is the adjacency indicator of the hybrid graph
defined in Section 2.4.2, and $c_{ij}$ is the pairwise structural
confidence defined in Section 2.4.2. The minimum
weight $q_{\min}$ prevents low-confidence regions from being completely
disconnected, whereas the exponent $\rho$ increases the relative contribution
of high-confidence contacts. The distance-decay term further reduces the
influence of more distant spatial neighbors.

Messages received by node $i$ are aggregated with a normalization floor:
\begin{equation}
    \mathbf m_i^{(\ell)} = \frac{\sum_j w_{ij}^{(\ell)}\, \mathbf m_{ij}^{(\ell)}}
    {\max\!\left(\sum_j w_{ij}^{(\ell)},\, 1\right)}.
    \label{eq:message_agg}
\end{equation}
The aggregated message is used to update the node representation with a
residual connection followed by layer normalization:
\begin{equation}
    \mathbf h_i^{(\ell+1)} = \operatorname{LN}\left\{\mathbf h_i^{(\ell)}
    + \phi_h^{(\ell)}\left[\mathbf h_i^{(\ell)} \parallel
    \mathbf m_i^{(\ell)}\right]\right\}.
    \label{eq:node_update}
\end{equation}

In addition to updating node features, the EGNN refines residue coordinates
using pairwise relative vectors. A bounded scalar coordinate-update
coefficient is generated from the edge message, and the coordinate update is
\begin{equation}
    \begin{aligned}
        \alpha_{ij}^{(\ell)} &= \tanh\left[\phi_x^{(\ell)}
        \left(\mathbf m_{ij}^{(\ell)}\right)\right] w_{ij}^{(\ell)},\\
        \mathbf x_i^{(\ell+1)} &= \mathbf x_i^{(\ell)} + 0.05\,
        \frac{\sum_j \alpha_{ij}^{(\ell)}
        \left(\mathbf x_i^{(\ell)} - \mathbf x_j^{(\ell)}\right)}
        {\max\!\left(\sum_j w_{ij}^{(\ell)},\, 1\right)}.
    \end{aligned}
    \label{eq:coord_update}
\end{equation}
The coordinates are recentered after every layer to reduce numerical drift.
Because the update depends only on scalar messages, distances, and relative
coordinate vectors, the operation remains equivariant to rigid rotations and
translations.

Three EGNN layers are stacked to progressively enlarge the spatial receptive
field. The first layer primarily encodes direct geometric neighbors, whereas
subsequent layers integrate higher-order spatial environments through
repeated message passing. Unlike approaches that simply append pLDDT as an
ordinary node feature, GHXTox uses pLDDT to directly modulate the strength of
geometric messages, allowing structural uncertainty to influence
representation learning at the propagation stage.

Continuous down-weighting does not imply that a low-pLDDT residue is
biologically unimportant. Low confidence may arise from prediction error or
potential conformational flexibility; the model reduces the reliability
assigned to uncertain coordinate evidence while preserving the residue's
sequence-semantic and physicochemical features.

\subsubsection{Full-Atom Chemical-Site Message Passing}
A C$\alpha$-based residue graph describes the global fold and spatial neighborhoods
but cannot distinguish chemically distinct regions within the same side
chain. For example, tyrosine contains both an aromatic ring and a hydroxyl
group, whereas tryptophan contains an aromatic system and an indole nitrogen.
Compressing each side chain into a single centroid would mix these chemically
different moieties. GHXTox therefore extracts chemistry-aware pseudo-sites
from the predicted full-atom PDB structure and processes them with a dedicated
local chemical branch.

Up to two chemical sites are retained for each residue. Site identity is
represented as an eight-dimensional multi-hot vector allowing overlapping
types:
\begin{equation}
    \mathbf t_{ir} \in \{0,1\}^{8},
    \label{eq:site_type}
\end{equation}
corresponding to positive, negative, donor, acceptor, aromatic, hydrophobic,
sulfur, and conformational, where $r$ indexes the chemical site within
residue $i$. Multiple types may be
assigned to the same site. For example, the cysteine sulfur site is labeled as
donor, hydrophobic, and sulfur-containing, whereas the proline ring is
labeled as hydrophobic and conformationally constrained. The aromatic ring
and hydroxyl group of tyrosine are represented as separate sites, as are the
aromatic system and the indole NE1 nitrogen (donor) of tryptophan.

A site coordinate is obtained from the mean position of the heavy atoms
defining the corresponding chemical moiety. Site orientation is represented by
either a bond direction or an unsigned plane normal. Because the sign of a
plane normal is arbitrary, downstream directional features use absolute
cosines, preserving rotation invariance and eliminating orientation-sign
ambiguity.

The initial representation of site $r$ in residue $i$ is
\begin{equation}
    \mathbf z_{ir}^{(0)} = \phi_{\mathrm{site}}\left[\mathbf t_{ir}
    \parallel q_i \parallel o_{ir}^{\mathrm{valid}}\right],
    \label{eq:site_init}
\end{equation}
where $q_i$ is the pLDDT score of the parent residue and
$o_{ir}^{\mathrm{valid}}$ indicates whether a reliable orientation could be
constructed; the input is projected into a 64-dimensional hidden space.

Chemical-site edges are established between valid sites separated by no more
than 16 Å. In the current configuration, same-residue site edges are
retained, preserving spatial relationships between chemically different
moieties in complex side chains. Each edge contains four complementary
feature families. First, the raw distance in Å is encoded using 16 radial
basis functions, $\boldsymbol\phi_{\mathrm{raw}}(d_{uv}) \in \mathbb{R}^{16}$.
Second, a peptide-size-normalized distance is constructed. The radius of
gyration of the C$\alpha$ coordinates is
\begin{equation}
    R_g = \sqrt{\frac{1}{L} \sum_{i=1}^{L}
    \left\lVert \mathbf x_i - \bar{\mathbf x} \right\rVert_2^2},
    \qquad \bar{\mathbf x} = \frac{1}{L} \sum_{i=1}^{L} \mathbf x_i,
    \qquad \bar{d}_{uv} = \frac{d_{uv}}{R_g},
    \label{eq:normalized_distance}
\end{equation}
where $\bar{\mathbf x}$ denotes the centroid of the valid C$\alpha$ coordinates, and
the normalized distance is encoded by eight radial basis functions over the
interval 0--4, $\boldsymbol\phi_{\mathrm{norm}}(\bar{d}_{uv}) \in
\mathbb{R}^{8}$. The raw representation describes the physical separation in
Å, whereas the normalized representation describes the separation relative to
the overall peptide size. Their combination allows the model to distinguish,
for example, an 8 Å contact in a compact short peptide from an 8 Å contact in
a larger peptide.

Third, four rotation-invariant orientation features are used:
\begin{equation}
    \mathbf o_{uv} = \left[|\mathbf o_u \cdot \mathbf r_{uv}|,\;
    |\mathbf o_v \cdot \mathbf r_{vu}|,\; |\mathbf o_u \cdot \mathbf o_v|,\;
    o_{uv}^{\mathrm{valid}}\right],
    \label{eq:orientation}
\end{equation}
which describe the alignment of the source orientation with the intersite
direction, the alignment of the target orientation with the reverse direction,
the mutual orientation of the two sites, and orientation availability.
Fourth, six explicit compatibility channels are defined:
\begin{equation}
    \mathbf c_{uv}^{\mathrm{chem}} = \left[c_{\mathrm{opp}}, c_{\mathrm{DA}},
    c_{\mathrm{aro}}, c_{\mathrm{cat}\text{-}\pi}, c_{\mathrm{SS}},
    c_{\mathrm{hyd}}\right],
    \label{eq:chem_channels}
\end{equation}
representing opposite-charge, donor--acceptor, aromatic--aromatic,
positive--aromatic (cation--$\pi$), sulfur--sulfur, and hydrophobic--hydrophobic
compatibility. These channels encode candidate chemical compatibility
patterns available to the model; they do not imply that a true salt bridge,
hydrogen bond, cation--$\pi$ interaction, disulfide bond, or hydrophobic core is
present in the predicted structure.

The full chemical-site edge representation is
\begin{equation}
    \mathbf e_{uv}^{\mathrm{chem}} = \left[\boldsymbol\phi_{\mathrm{raw}}
    (d_{uv}) \parallel \boldsymbol\phi_{\mathrm{norm}}(\bar{d}_{uv})
    \parallel \mathbf o_{uv} \parallel \mathbf c_{uv}^{\mathrm{chem}}
    \parallel \sqrt{q_u q_v}\right].
    \label{eq:site_edge}
\end{equation}
At layer $l$, the site message, its pairwise-confidence-weighted form, and the
confidence-normalized aggregation are
\begin{equation}
    \begin{aligned}
        \mathbf m_{uv}^{(l)} &= \phi_{\mathrm{chem}}^{(l)}
        \left[\mathbf z_u^{(l)} \parallel \mathbf z_v^{(l)}
        \parallel \mathbf e_{uv}^{\mathrm{chem}}\right],\\
        \tilde{\mathbf m}_{uv}^{(l)} &= \sqrt{q_u q_v}\,
        \mathbf m_{uv}^{(l)},\\
        \bar{\mathbf m}_v^{(l)} &= \frac{\sum_{u \in \mathcal{N}(v)}
        \tilde{\mathbf m}_{uv}^{(l)}}
        {\max\!\left(\sum_{u \in \mathcal{N}(v)} \sqrt{q_u q_v},\, 1\right)}.
    \end{aligned}
    \label{eq:site_message}
\end{equation}
After two message-passing layers, valid sites belonging to the same residue
are mean-pooled, and a residue-level chemical residual is obtained:
\begin{equation}
    \begin{aligned}
        \Delta \mathbf h_i^{\mathrm{chem}} &= \mathbf W_{\mathrm{chem}}\,
        \underset{r \in \mathcal{S}_i}{\operatorname{Mean}}
        \left(\mathbf z_{ir}^{(2)}\right),\\
        \mathbf W_{\mathrm{chem}}^{(0)} &= 0, \quad
        \mathbf b_{\mathrm{chem}}^{(0)} = 0,\\
        \mathbf h_i^{\mathrm{sp}} &= \mathbf h_i^{\mathrm{EGNN}} +
        \Delta \mathbf h_i^{\mathrm{chem}},
    \end{aligned}
    \label{eq:chem_residual}
\end{equation}
where the projection $\mathbf W_{\mathrm{chem}}$ and its bias
$\mathbf b_{\mathrm{chem}}$ are initialized to zero, so that
$\Delta \mathbf h_i^{\mathrm{chem}} = 0$ at the beginning of training. This
conservative residual design allows the new chemical-site branch to
progressively learn useful corrections from an intact geometric backbone
rather than perturbing the existing spatial representation with a randomly
initialized contribution.

\subsubsection{Availability-Masked Conformer-Ensemble Residual}
A single predicted structure provides a deterministic spatial anchor, whereas
a peptide may adopt multiple related conformations. Running the complete
geometric network independently on every conformer would be computationally
expensive and could allow the model to exploit differences between conformer
generators or uneven ensemble coverage. The current GHXTox model therefore
uses a lightweight conformer-statistics residual. The single-structure expert
remains the base predictor, and ensemble statistics are permitted to correct
its logit only when valid conformer information is available.

For a peptide with $M$ aligned conformers, the residue-level
root-mean-square fluctuation is
\begin{equation}
    \mathrm{RMSF}_i = \sqrt{\frac{1}{M} \sum_{m=1}^{M}
    \left\lVert \mathbf x_i^{(m)} - \bar{\mathbf x}_i \right\rVert_2^2},
    \qquad \bar{\mathbf x}_i = \frac{1}{M} \sum_{m=1}^{M}
    \mathbf x_i^{(m)}.
    \label{eq:rmsf}
\end{equation}
For each residue pair $i, j$, the mean distance and the distance standard
deviation are
\begin{equation}
    \bar{d}_{ij} = \frac{1}{M} \sum_{m=1}^{M} d_{ij}^{(m)}, \qquad
    \sigma_{ij}^{d} = \sqrt{\frac{1}{M} \sum_{m=1}^{M}
    \left(d_{ij}^{(m)} - \bar{d}_{ij}\right)^2}.
    \label{eq:pair_stats}
\end{equation}
For a distance cutoff $r \in \{6, 8, 10\}$ Å, the contact occupancy and the
8 Å contact-occupancy uncertainty are
\begin{equation}
    O_{ij}^{(r)} = \frac{1}{M} \sum_{m=1}^{M}
    \mathbb{I}\left(d_{ij}^{(m)} < r\right), \qquad
    U_{ij}^{(8)} = O_{ij}^{(8)}\left(1 - O_{ij}^{(8)}\right).
    \label{eq:occupancy}
\end{equation}
Values close to zero or one indicate contacts that are consistently absent or
present across the sampled conformers, whereas values near 0.5 indicate
contact switching.

The ensemble statistics are organized into an eight-dimensional residue-level
descriptor $\mathbf r_i^{\mathrm{conf}} \in \mathbb{R}^{8}$, containing the
residue RMSF; the mean pairwise distance to other residues; the standard
deviation of the mean pairwise distances associated with the residue; the
mean pairwise-distance standard deviation; the mean contact occupancies at 6,
8, and 10 Å; and the mean 8 Å contact-occupancy uncertainty, which together
form the eight-dimensional descriptor. A twelve-dimensional peptide-level
descriptor $\mathbf g^{\mathrm{conf}} \in \mathbb{R}^{12}$ is also
constructed, containing the normalized number of conformers; the mean,
standard deviation, and maximum of RMSF; the mean, standard deviation, and
maximum of the conformer RMSD; the mean and standard deviation of the radius
of gyration; the mean and maximum of the pair-distance standard deviation;
and the global mean of $U_{ij}^{(8)}$, which together form the
twelve-dimensional descriptor. Before neural encoding, distance- and
fluctuation-based statistics are rescaled by fixed constants to place
heterogeneous descriptors on comparable numerical ranges; the exact scaling
factors are reported in the implementation-details table.

The residue-level ensemble features are encoded and aggregated using both
masked mean and masked maximum pooling, whereas the global descriptor is
processed by a separate encoder:
\begin{equation}
    \mathbf h_{\mathrm{conf}}^{\mathrm{res}} = \left[\mathrm{MeanPool}
    \left(\phi_r\left(\mathbf R^{\mathrm{conf}}\right)\right) \parallel
    \mathrm{MaxPool}\left(\phi_r\left(\mathbf R^{\mathrm{conf}}\right)\right)
    \right], \qquad
    \mathbf h_{\mathrm{conf}}^{\mathrm{global}} = \phi_g
    \left(\mathbf g^{\mathrm{conf}}\right).
    \label{eq:conf_embeddings}
\end{equation}
The generator is encoded by a three-dimensional one-hot vector
$\mathbf s_{\mathrm{conf}}\in\mathbb{R}^{3}$ for PepFlow,
AlphaFold2/ColabFold, or another admitted generator. A learned source
embedding $\phi_s(\mathbf s_{\mathrm{conf}})$ is concatenated with the base
structural embedding and the two ensemble embeddings to produce a raw logit
residual:
\begin{equation}
    \Delta z_{\mathrm{conf}}^{\mathrm{raw}} = \phi_{\mathrm{conf}}\left[
    \mathbf h_{\mathrm{base}} \parallel \mathbf h_{\mathrm{conf}}^{\mathrm{res}}
    \parallel \mathbf h_{\mathrm{conf}}^{\mathrm{global}}
    \parallel \phi_s(\mathbf s_{\mathrm{conf}})\right].
    \label{eq:raw_residual}
\end{equation}
The final linear layer of the conformer residual head is initialized to zero.
Using the availability indicator $c_{\mathrm{conf}}$ defined in
Eq.~\eqref{eq:conf_available}, the learnable conformer weight, the effective
residual, and the final structural-expert logit are
\begin{equation}
    \begin{aligned}
        \beta_{\mathrm{conf}}^{(0)} &= \mathrm{logit}(0.25),\\
        \lambda_{\mathrm{conf}} &= \sigma\left(\beta_{\mathrm{conf}}\right),\\
        \Delta z_{\mathrm{conf}} &= c_{\mathrm{conf}}\, \lambda_{\mathrm{conf}}\,
        \Delta z_{\mathrm{conf}}^{\mathrm{raw}},\\
        z_{\mathrm{SE3D}} &= z_{\mathrm{base}} + \Delta z_{\mathrm{conf}},
        \qquad c_{\mathrm{conf}} = 0 \;\Rightarrow\; z_{\mathrm{SE3D}} = z_{\mathrm{base}}.
    \end{aligned}
    \label{eq:effective_residual}
\end{equation}
where the effective conformer weight is initialized to 0.25 at the
beginning of training through $\beta_{\mathrm{conf}}^{(0)} =
\mathrm{logit}(0.25)$, so that $\sigma(\beta_{\mathrm{conf}}^{(0)}) = 0.25$. Zero-filled feature tensors alone would not guarantee a zero
correction because the residual network may contain nonzero biases; the
additional hard availability mask ensures an exact fallback to the
single-structure prediction. The conformer branch is permitted only to refine
the single-structure predictor rather than replace the geometric backbone,
reducing the risk of overfitting to limited and uneven ensemble coverage.

\subsubsection{pLDDT-Aware Sequence--Structure Cross-Attention}
After the EGNN and the chemical-site residual are applied, the model obtains
the chemical-enhanced spatial representation
\begin{equation}
    \mathbf H^{\mathrm{sp}} = \left[\mathbf h_1^{\mathrm{sp}}, \ldots,
    \mathbf h_L^{\mathrm{sp}}\right]^{\top} \in \mathbb{R}^{L \times 128},
    \label{eq:spatial_enhanced}
\end{equation}
where $\mathbf h_i^{\mathrm{sp}}$ is defined in Eq.~\eqref{eq:chem_residual}.
Spatial information alone does not determine whether a structural pattern is
relevant to the current sequence context. For example, similar hydrophobic
contacts may have different implications depending on residue identities and
sequence positions. GHXTox therefore uses pLDDT-aware cross-attention within
the structural expert, allowing contextualized sequence states to selectively
retrieve spatial evidence.

A node-level gate is first constructed for each residue. The base gate uses
the 42-dimensional residue descriptors and the pLDDT score:
\begin{equation}
    \begin{aligned}
        g_i^{\mathrm{base}} &= \phi_{\mathrm{node}}\left[\mathbf f_i^{\mathrm{res}} \parallel q_i\right],\\
        \mathbf c_i^{\mathrm{node}} &= \left[\frac{q_i - \mu_q}{T_q},\; \frac{(q_i - \mu_q)^2}{T_q^2},\; \mathbb{I}(q_i < 0.55)\right],\\
        \mu_q &= 0.70, \quad T_q = 0.12,
    \end{aligned}
    \label{eq:node_gate_base}
\end{equation}
where the residual confidence vector $\mathbf c_i^{\mathrm{node}}$ increases
the sensitivity of the gate to the distribution of structural confidence. The
residual scale and the final node gate are
\begin{equation}
    r_i^{\mathrm{node}} = 0.5 + \sigma\left(\mathbf w_n^{\top}
    \mathbf c_i^{\mathrm{node}} + b_n\right), \qquad
    g_i^{\mathrm{node}} = \mathrm{clamp}\left(g_i^{\mathrm{base}}\,
    r_i^{\mathrm{node}},\; 0, 1\right),
    \label{eq:node_gate_final}
\end{equation}
and the spatial representation is continuously down-weighted as
\begin{equation}
    \widehat{\mathbf H}^{\mathrm{sp}} = \mathbf H^{\mathrm{sp}} \odot
    \mathbf g^{\mathrm{node}}, \qquad
    \mathbf g^{\mathrm{node}} = \left[g_1^{\mathrm{node}}, \ldots,
    g_L^{\mathrm{node}}\right]^{\top} \in [0,1]^L,
    \label{eq:gated_spatial}
\end{equation}
where the node-level gate vector $\mathbf g^{\mathrm{node}}$ is broadcast
across the feature dimension and each gate value is clamped to $[0,1]$.

In the cross-attention operation, the contextualized sequence representation
serves as the query, and the gated spatial representation serves as the key
and value:
\begin{equation}
    \mathbf H^{\mathrm{cross}} = \mathrm{MHA}\left(Q = \mathbf H^{\mathrm{seq}},
    K = \widehat{\mathbf H}^{\mathrm{sp}},
    V = \widehat{\mathbf H}^{\mathrm{sp}}\right).
    \label{eq:cross_attn_sp}
\end{equation}
This direction has a direct interpretation: each residue position retrieves
spatial evidence according to its contextual sequence representation. The
operation is not a direct feature interaction between the LoRA-ProtT5--Mamba
expert and the structural expert; it occurs internally between the
ESM2-enhanced sequence states and the geometric states of the structural
expert.

In addition to residue-level gating, a peptide-level gate is generated from
the overall structure-confidence profile. The model first computes
\begin{equation}
    \bar{q} = \frac{1}{L}\sum_{i=1}^{L} q_i, \qquad
    q_{\min} = \min_i q_i, \qquad
    r_{\mathrm{low}} = \frac{1}{L}\sum_{i=1}^{L}
    \mathbb{I}(q_i < 0.55),
    \label{eq:peptide_stats}
\end{equation}
and constructs the peptide-level confidence vector
$\mathbf c^{\mathrm{pep}} = \left[(\bar{q} - \mu_q)/T_q,\;
(q_{\min} - \mu_q)/T_q,\; r_{\mathrm{low}}\right]$. The base peptide-level
gate and its residual-adjusted form are
\begin{equation}
    \begin{aligned}
        g_{\mathrm{pep}}^{\mathrm{base}} &= \phi_{\mathrm{global}}\left[\bar{q},\; q_{\min},\; r_{\mathrm{low}}\right],\\
        r_{\mathrm{pep}} &= 0.5 + \sigma\left(\mathbf w_{\mathrm{pep}}^{\top} \mathbf c^{\mathrm{pep}} + b_{\mathrm{pep}}\right),\\
        g_{\mathrm{pep}} &= \mathrm{clamp}\left(g_{\mathrm{pep}}^{\mathrm{base}}\, r_{\mathrm{pep}},\; 0, 1\right),
    \end{aligned}
    \label{eq:peptide_gate}
\end{equation}
where the peptide-level residual scale is analogous to, but separately
parameterized from, the node-level residual scale, and $g_{\mathrm{pep}} \in
[0,1]$ denotes the scalar peptide-level structural gate.

The final residue-level fusion representation is
\begin{equation}
    \mathbf H^{\mathrm{fused}} = \operatorname{LN}\left[\mathbf H^{\mathrm{seq}}
    \parallel g_{\mathrm{pep}}\, \mathbf H^{\mathrm{cross}}\right].
    \label{eq:fusion_repr}
\end{equation}
The two gating levels answer complementary questions: how reliable the
structural evidence is at an individual residue position, and how strongly the
structural branch should be trusted for the peptide as a whole. The sequence
representation is always retained, whereas the structural context is
dynamically scaled according to prediction confidence. Consequently, when the
predicted structure has low overall confidence, the structural expert can
still rely on its internal sequence-semantic pathway without allowing
uncertain coordinates to dominate the classification.

\subsubsection{Global Feature Integration and Structural Classification}
After pLDDT-aware cross-attention, each valid residue is represented by a
256-dimensional fused vector obtained by concatenating the 128-dimensional
internal sequence state with the 128-dimensional structural context:
\begin{equation}
    \mathbf H^{\mathrm{fused}} \in \mathbb{R}^{L \times 2 d_h},
    \qquad d_h = 128,
    \label{eq:fused_repr}
\end{equation}
so that $\mathbf H^{\mathrm{fused}} \in \mathbb{R}^{L \times 256}$. Masked
mean pooling is applied over the valid residue positions:
\begin{equation}
    \mathbf h_{\mathrm{pool}} = \frac{\sum_{i=1}^{L} m_i\,
    \mathbf H^{\mathrm{fused}}_i}{\sum_{i=1}^{L} m_i} \in \mathbb{R}^{256},
    \label{eq:masked_pool_3d}
\end{equation}
where the denominator equals the number of valid residue positions and the
mask excludes padding positions from the peptide representation. In contrast
to maximum pooling, mean pooling integrates structural evidence distributed
across multiple residue positions, which is appropriate when
toxicity-related spatial patterns are not confined to a single local region.

The residue graph primarily models local and intermediate-range environments,
whereas a finite number of message-passing layers may not explicitly recover
all peptide-level physicochemical statistics. GHXTox therefore includes five
global physicochemical descriptors, normalized to comparable scales before
concatenation:
\begin{equation}
    \mathbf f_{\mathrm{pep}}^{\mathrm{phys}} = \left[\frac{L}{128},\;
    \frac{Q_{\mathrm{net}}}{L},\; f_{\mathrm{aro}},\; f_{\mathrm{Cys}},\;
    \frac{\bar{h}_{\mathrm{hyd}}}{4.5}\right] \in \mathbb{R}^{5},
    \label{eq:global_descriptors}
\end{equation}
representing the normalized peptide length, the length-normalized net charge,
the aromatic-residue fraction, the cysteine fraction, and the scaled mean
hydropathy. The classifier input is
\begin{equation}
    \mathbf h_{\mathrm{cls}} = \left[\mathbf h_{\mathrm{pool}} \parallel
    \mathbf f_{\mathrm{pep}}^{\mathrm{phys}}\right] \in \mathbb{R}^{261}.
    \label{eq:classifier_input}
\end{equation}
The base structural logit is produced by a three-layer multilayer perceptron,
denoted compactly as
\begin{equation}
    z_{\mathrm{base}} = \phi_{\mathrm{cls}}\left(\mathbf h_{\mathrm{cls}}\right),
    \label{eq:mlp_logit}
\end{equation}
with dimensions $261 \rightarrow 128 \rightarrow 64 \rightarrow 1$, where each
hidden layer applies GELU activation and dropout.

Applying the availability-masked conformer correction defined in
Eq.~\eqref{eq:effective_residual} yields the final structural-expert logit,
and the structural-expert probability is
\[
    z_{\mathrm{SE3D}} = z_{\mathrm{base}} + c_{\mathrm{conf}}\,
    \lambda_{\mathrm{conf}}\, \Delta z_{\mathrm{conf}}^{\mathrm{raw}},
    \qquad p_{\mathrm{SE3D}} = \sigma\left(z_{\mathrm{SE3D}}\right),
\]
as defined in Eq.~\eqref{eq:effective_residual}. No standalone structural
threshold is used to generate the final GHXTox class at this stage. Instead,
$p_{\mathrm{SE3D}}$ is passed directly to the top-level dual-expert fusion
described in Section 2.2.3, where it is combined with the calibrated 1D
expert probability $p_{\mathrm{1D}}$ using frozen fusion weights. A standalone
structural threshold is used only for structural-branch ablations or
independent expert evaluation and does not determine the final GHXTox
prediction.

The structural expert can therefore be summarized as
\[
    \begin{aligned}
        &\text{sequence-enhanced residue semantics} + \text{confidence-aware
        geometry}\\
        &+ \text{full-atom chemical sites} + \text{global physicochemical
        descriptors}\\
        &+ \text{availability-masked conformer statistics}
        \longrightarrow p_{\mathrm{SE3D}}.
    \end{aligned}
\]
The central contribution is not merely the inclusion of additional structural
features, but the explicit separation of structural information by spatial
scale and reliability. The geometric graph captures global residue
relationships, the chemical-site branch resolves local functional moieties,
the conformer residual summarizes sampled structural variation, and pLDDT
continuously modulates edge propagation, chemical messages, and
sequence--structure fusion. This design allows the structural expert to
incorporate complementary 3D evidence while handling structure-prediction
uncertainty in an explicit and auditable manner.

Overall, the structural expert integrates residue-level sequence semantics,
confidence-aware geometry, full-atom chemical-site interactions,
peptide-level physicochemical properties, and availability-masked conformer
statistics. Rather than treating all predicted structural evidence as equally
reliable, GHXTox continuously modulates its contribution using residue-level
and peptide-level confidence signals. The following section describes the
training, out-of-fold model-selection, calibration, and parameter-freezing
protocols used to optimize the complete framework.

\subsection{Training and Model-Selection Protocol}
GHXTox consists of a task-adaptive 1D sequence expert and a
sequence-enhanced 3D structural expert. Because the two experts differ
substantially in input representation, parameter scale, and optimization
behavior, they are not jointly optimized through a single end-to-end
objective. Instead, expert-level representations are learned independently,
after which probability calibration and top-level fusion are performed using
training out-of-fold (OOF) predictions. This strategy allows each expert to be
optimized under modality-appropriate settings while separating representation
learning, checkpoint selection, probability calibration, fusion-parameter
determination, and benchmark evaluation into clearly defined stages.

The training protocol follows three central principles. First, identical or
highly similar peptides must not appear simultaneously in a training
partition and the corresponding validation partition. Second, every
prediction used for calibration, fusion-weight selection, or threshold
determination must be generated by a model that was not trained on the
corresponding sample. Third, Test1 and Test2 are not used for checkpoint
selection, calibrator fitting, fusion search, or decision-threshold
optimization in the current pipeline. Their prior observation during project
development is disclosed, and they are therefore treated as frozen
retrospective benchmarks rather than newly blinded external validation sets.

\subsubsection{Group-Aware Five-Fold OOF Generation}
The training set contains $N = 6387$ peptides. To reduce performance
inflation caused by highly similar sequences, a fixed five-fold manifest is
constructed by clustering the training sequences into homology groups. The fold-level homology groups were
constructed through a separate sequence-clustering procedure, distinct from
the dataset-level redundancy removal described in Section 2.1. The first operation determines which samples remain in the
benchmark dataset, whereas the second determines which retained samples must
remain in the same cross-validation fold. Let
\begin{equation}
    \mathcal{D} = \left\{\left(\mathcal{S}_i, y_i, g_i\right)\right\}_{i=1}^{N},
    \qquad N = 6387,
    \label{eq:train_data}
\end{equation}
denote the training data, where $\mathcal{S}_i$ is the peptide sequence,
$y_i \in \{0,1\}$ is its toxicity label, and $g_i$ denotes its
sequence-homology group. The data are partitioned into five disjoint
validation folds:
\begin{equation}
    \begin{aligned}
        \mathcal{D} &= \bigcup_{k=1}^{5}\mathcal{D}_k, &
        \mathcal{D}_k \cap \mathcal{D}_l &= \varnothing, \quad k \neq l,\\
        g_i = g_j &\Longrightarrow f(i) = f(j),
    \end{aligned}
    \label{eq:folds}
\end{equation}
where $f(i)$ denotes the fold assigned to sample $i$. All peptides belonging
to the same homology group are therefore placed in the same fold, preventing
closely related sequences from crossing the training--validation boundary.

For fold $k$, the model is trained only on the complement
$\mathcal{D}_{-k} = \mathcal{D} \setminus \mathcal{D}_k$ and evaluated on the
unseen fold $\mathcal{D}_k$. The OOF probability of sample $i \in
\mathcal{D}_k$ is defined as
\begin{equation}
    \mathcal{D}_{-k} = \mathcal{D} \setminus \mathcal{D}_k, \qquad
    \widehat{p}_i^{\mathrm{OOF}} = F_{\boldsymbol\theta^{(-k)}}\left(
    \mathcal{S}_i\right), \quad i \in \mathcal{D}_k,
    \label{eq:oof_prob}
\end{equation}
where $\boldsymbol\theta^{(-k)}$ denotes the parameters learned exclusively
from $\mathcal{D}_{-k}$. Each training peptide consequently receives exactly
one probability from a model that did not observe that peptide during
optimization.

For the 1D expert, full seed-specific OOF predictions are generated for the
training set, and the three probability vectors are aligned by sample index
and averaged:
\begin{equation}
    \overline{p}_{i,\mathrm{1D}}^{\mathrm{OOF}} = \frac{1}{3}
    \sum_{s \in \{42, 123, 2025\}}
    \widehat{p}_{i,\mathrm{1D}}^{(s,\mathrm{OOF})}.
    \label{eq:seed_ensemble}
\end{equation}
The resulting OOF ensemble scores are used to fit the Platt calibrator
described in Section 2.3.4.

The structural expert uses the same fixed group-aware folds. Each
fold-specific checkpoint produces probabilities only for its held-out
validation fold:
\begin{equation}
    \widehat{p}_{i,\mathrm{SE3D}}^{\mathrm{OOF}} = G_{\boldsymbol\phi^{(-k)}}
    \left(\mathcal{S}_i, \mathcal{X}_i, \mathcal{A}_i, \mathbf q_i,
    \mathcal{C}_i\right), \qquad i \in \mathcal{D}_k,
    \label{eq:oof_se3d}
\end{equation}
where $\mathcal{X}_i$, $\mathcal{A}_i$, $\mathbf q_i$, and $\mathcal{C}_i$
denote the residue geometry, full-atom chemical information, pLDDT
confidence, and available conformer-ensemble statistics, respectively;
$\mathcal{C}_i$ may be unavailable and is handled by the conformer-availability
mask defined in Section 2.4.1.

Before combining the expert-level OOF predictions, the implementation
verifies immutable sample indices (\texttt{source\_index}), peptide
sequences, toxicity labels, and fold assignments. Duplicate sample indices or
mismatches in sequence, label, or fold assignment cause the fusion procedure
to terminate rather than silently aligning prediction files by row order. OOF
predictions are treated as \textbf{training-only model-selection assets}.
Although OOF predictions are generated by models that did not observe the
corresponding validation samples, parameters selected from the complete OOF
table remain development-stage choices and must be evaluated on independent
test sets.

\subsubsection{Loss Functions and Class-Imbalance Handling}
The 1D and 3D experts are optimized using separate objectives. This prevents
uneven structural coverage, conformer availability, and confidence variation
from directly perturbing the language-model optimization, and it avoids the
large ProtT5 parameter space from dominating the smaller structure-specific
residual modules.

\textbf{Loss for the 1D expert.}
As described in Section 2.3.3, the sequence expert is trained using
positive-class-weighted binary cross-entropy:
\begin{equation}
    \mathcal{L}_{\mathrm{1D}} = -\frac{1}{B} \sum_{i=1}^{B}
    \left[w_+ y_i \log \sigma(z_i) + \left(1 - y_i\right)
    \log\left(1 - \sigma(z_i)\right)\right], \qquad
    w_+ = \frac{N_-}{N_+},
    \label{eq:loss_1d}
\end{equation}
where the class counts $N_+$ and $N_-$ are independently calculated from each
fold's training partition. The current 1D training pipeline does not use
Borderline-SMOTE or sequence point-mutation augmentation. The weighted loss
increases the contribution of the minority toxic class without introducing
synthetic near-duplicate peptides.

\textbf{Base geometric and chemical-site stage.}
The structural expert is optimized in stages. The full-atom chemical-site
branch is first initialized from a previously trained pLDDT-aware geometric
model, while the remaining backbone is frozen. This stage uses
positive-class-weighted BCE:
\begin{equation}
    \mathcal{L}_{\mathrm{chem}} = -\frac{1}{B} \sum_{i=1}^{B}
    \left[w_+ y_i \log \sigma\left(z_i^{\mathrm{base}}\right) +
    \left(1 - y_i\right) \log\left(1 -
    \sigma\left(z_i^{\mathrm{base}}\right)\right)\right],
    \label{eq:loss_chem}
\end{equation}
where $w_+$ is defined as in Section 2.3.3 and is recalculated from each
fold's training partition.

During training, pLDDT-dependent Gaussian coordinate perturbations are
applied to reduce sensitivity to local errors in predicted structures. The
perturbation scale for residue $i$ is
\begin{equation}
    \sigma_i = \sigma_0\left[1 + \kappa\left(1 - q_i\right)\right],
    \qquad \sigma_0 = 0.12\ \text{\AA}, \quad \kappa = 2.0,
    \label{eq:perturb_scale}
\end{equation}
and residue-wise perturbations are drawn as
\begin{equation}
    \boldsymbol\varepsilon_i \sim \mathcal{N}\left(\mathbf 0,
    \sigma_i^2 \mathbf I_3\right), \qquad
    \widetilde{\mathbf x}_i = \mathbf x_i + \boldsymbol\varepsilon_i.
    \label{eq:perturb}
\end{equation}
The same residue-wise displacement is applied consistently to the
corresponding backbone, functional-group, and chemical-site coordinates,
thereby preserving their intra-residue relative geometry. No coordinate
perturbation is applied during validation or inference. Low-pLDDT residues
therefore receive stronger perturbations than high-confidence residues. This
augmentation uses pLDDT to modulate the perturbation magnitude but does not
treat pLDDT as an exact estimate of coordinate-error variance.

\textbf{Conformer-ensemble residual stage.}
During conformer-residual training, the pLDDT-aware geometric backbone,
internal sequence encoder, chemical-site branch, and base classifier are
frozen. Only the residue-level conformer encoder, the global conformer
encoder, the conformer-residual MLP, and the conformer logit-residual
parameters are optimized. This stage is initialized from the trained
chemical-site checkpoint. To avoid applying a second class-frequency
correction on top of the existing base classifier, the current
conformer-residual configuration uses unweighted BCE and disables coordinate
perturbation, with a residual-magnitude penalty:
\begin{equation}
    \mathcal{L}_{\mathrm{conf}} = \operatorname{BCEWithLogits}\left(z_{\mathrm{SE3D}},
    y\right) + \lambda_\Delta\,
    \frac{\sum_{i=1}^{B} c_i \left|\Delta z_{\mathrm{conf},i}\right|}
    {\max\!\left(\sum_{i=1}^{B} c_i,\, 1\right)},
    \qquad \lambda_\Delta = 0.005,
    \label{eq:loss_conf}
\end{equation}
where the residual-magnitude term is computed only on samples with available
conformer information. It limits large logit corrections and encourages the
conformer branch to act as a conservative refinement rather than re-learning
the complete toxicity decision boundary.

The final conformer-residual layer is zero-initialized, and its output is
additionally controlled by the hard \texttt{conformer\_available} mask. Zero
initialization makes the initial model equivalent to the single-structure
expert, whereas the hard mask guarantees a zero residual for samples without
conformer evidence. Together, these mechanisms provide a stable incremental
extension of the structural model.

\subsubsection{Optimization, Checkpoint Selection, and Hyperparameter Settings}
Both experts are optimized using AdamW \cite{loshchilov2019adamw},
but different learning rates, training durations, and checkpoint-selection
criteria are employed. All random-number generators are initialized with
fixed seeds. Data are shuffled only during training, whereas validation and
inference loaders preserve deterministic ordering.
Table~\ref{tab:training_settings} summarizes the optimization and
model-selection settings of both experts.

\begin{table}[htbp]
    \centering
    \caption{Main optimization and model-selection settings of the two GHXTox
             experts.}
    \label{tab:training_settings}
    \resizebox{\textwidth}{!}{\begin{tabular}{lccccp{2.9cm}}
        \toprule
        Module or stage & Batch size & Learning rate & Weight decay &
        Epochs & Checkpoint rule \\
        \midrule
        1D LoRA-ProtT5 + Bi-Mamba & 16 & $3\times10^{-4}$ &
        $1\times10^{-2}$ & Fixed 15 & No early stopping \\
        3D chemical-site refinement & 32 & $5\times10^{-4}$ &
        $1\times10^{-4}$ & Max. 30 & Best validation MCC; patience 8 \\
        Conformer residual refinement & 32 & $1\times10^{-3}$ &
        $5\times10^{-4}$ & Max. 50 & Minimum validation loss; patience 8 \\
        Single-structure chemical-site backbone & 32 & Stage-specific & Stage-specific &
        Fixed 8 & Median fold-selected epoch \\
        Source-aware full2795 conformer residual & 32 & $1\times10^{-3}$ &
        $5\times10^{-4}$ & Fixed 6 & OOF-defined refit protocol \\
        \bottomrule
    \end{tabular}}
\end{table}

\textbf{1D expert.}
The LoRA-ProtT5 and bidirectional Mamba expert is trained with a batch size
of 16, a learning rate of $3\times10^{-4}$, a weight decay of
$1\times10^{-2}$, a fixed duration of 15 epochs, dropout of 0.1, FP16 mixed
precision, and ProtT5 gradient checkpointing, using random seeds 42, 123, and
2025. This configuration yields approximately 13.1 million trainable
parameters, comprising the LoRA updates and the downstream
sequence-classification modules, whereas the pretrained base transformations
and embeddings remain frozen. Every fold model and every full-data model is
trained for a fixed 15 epochs. Validation-based early stopping and a separate
learning-rate scheduler are not used. A fixed training budget ensures
consistency across folds and random seeds and avoids differences in training
duration caused by noisy fold-level validation metrics. After
cross-validation, three models are independently trained on the complete
training set using the same configuration. The deployed sequence probability
is obtained by averaging these three full-data models rather than selecting
the seed with the highest validation performance.

\textbf{3D chemical-site stage.}
The full-atom chemical-site refinement uses a batch size of 32, a learning
rate of $5\times10^{-4}$, a weight decay of $1\times10^{-4}$, a maximum
duration of 30 epochs, early-stopping patience of 8 epochs, base structural
dropout of 0.22, and chemical-site dropout of 0.1. This stage is initialized
from a pLDDT-aware geometric checkpoint, and only the chemical-site branch is
made trainable. Training is terminated when the validation MCC does not
improve for eight consecutive epochs, and the checkpoint with the highest
validation MCC is retained.

\textbf{Conformer-residual stage.}
The conformer-ensemble residual stage uses a batch size of 32, a learning
rate of $1\times10^{-3}$, a weight decay of $5\times10^{-4}$, a maximum
duration of 50 epochs, early-stopping patience of 8 epochs,
conformer-residual dropout of 0.25, a conformer hidden dimension of 32, an
initial conformer weight of 0.25, and a residual L1 coefficient of 0.005.
Validation loss is used instead of training accuracy to select the conformer
checkpoint, reducing sensitivity to threshold fluctuations arising from
limited conformer coverage. Only conformer-related modules are optimized,
whereas the geometric and chemical backbone remains frozen.

During structural-expert training, gradients of all trainable parameters are
clipped to a maximum norm of 5.0
($\lVert\nabla_{\phi}\mathcal{L}\rVert_2 \leq 5.0$), limiting unusually
large gradients caused by local residual corrections or coordinate
perturbations.

When a full-data structural checkpoint is required, the number of refitting
epochs is not selected using either benchmark test set. Instead, a fixed
epoch count is obtained from the median of the fold-specific selected epochs.
For the single-structure chemical-site backbone, the recorded fold-level
epochs are 8, 8, 14, 2, and 4, whose median is 8, and the full-data refit is
trained for eight fixed epochs. The retained source-aware conformer residual
uses the full2795 cohort and a six-epoch full-data refit fixed from the
training-side protocol. Neither the refit duration nor the checkpoint was
selected using Test1 or Test2.

% TODO: report the frozen hardware/software environment (GPU model, CUDA,
%       PyTorch, Transformers, training time) from the final reproducibility
%       manifest; do not guess before the environment is frozen.

\subsubsection{OOF-Based Calibration, Fusion-Weight Selection, and Threshold Freezing}
After the expert-level OOF predictions have been generated, model selection
proceeds at the probability level. No ProtT5, Mamba, EGNN, chemical-site, or
conformer-residual parameters are updated during this stage. Only the training
OOF probabilities are used to determine the calibrator, fusion weight, and
final decision threshold.

The three seed-specific 1D OOF probabilities are first averaged and
calibrated using the monotonic Platt mapping described in Section 2.3.4:
\begin{equation}
    p_{i,\mathrm{1D}}^{\mathrm{cal}} = \sigma\left[a\,
    \mathrm{logit}\left(\bar{p}_{i,\mathrm{1D}}^{\mathrm{OOF}}\right) + b
    \right], \qquad a = \exp(\eta) > 0,
    \label{eq:platt_oof}
\end{equation}
where the scale is parameterized as a positive value, $a = \exp(\eta) > 0$, as
described in Section 2.3.4. Platt calibration is applied only to the 1D
ensemble. The structural-expert probability is not subjected to an additional
top-level calibrator, and no new calibrator is fitted after the final fusion.

The calibrated sequence probability and the structural-expert probability are
combined in probability space:
\begin{equation}
    p_i(w) = w\, p_{i,\mathrm{1D}}^{\mathrm{cal}} + \left(1 - w\right)\,
    p_{i,\mathrm{SE3D}}^{\mathrm{OOF}}, \qquad 0 \leq w \leq 1,
    \label{eq:fusion_prob}
\end{equation}
where $w$ denotes the 1D expert weight and $1 - w$ denotes the
sequence-enhanced structural-expert weight. No additional stacking classifier
or sample-specific routing network is fitted at this stage, reducing the risk
of second-level overfitting to the finite OOF table.

For this top-level fusion, $p_{i,\mathrm{SE3D}}^{\mathrm{OOF}}$ is the pure
\texttt{structure\_probability} emitted by the pLDDT-aware structural expert.
It is not the probability of an earlier sequence--structure fusion model and
contains no contribution from the current or historical 1D predictor. This
strict interface prevents duplicated sequence evidence and avoids mixing old
and new model generations during OOF weight selection.

For a candidate weight $w$ and threshold $\tau$, the predicted class is
\begin{equation}
    \hat{y}_i(w, \tau) = \mathbb{I}\left[p_i(w) \geq \tau\right],
    \label{eq:decision_rule}
\end{equation}
and the frozen parameters are selected according to
\begin{equation}
    (w^*, \tau^*) = \arg\max_{w,\tau}\;
    \mathrm{MCC}\left(\mathbf y, \hat{\mathbf y}(w, \tau)\right).
    \label{eq:selection}
\end{equation}
MCC is used as the primary objective because of the pronounced class
imbalance in the benchmark data, particularly in Test2. Balanced accuracy was
used as the secondary selection criterion when candidate operating points
produced identical MCC values under the evaluated search procedure.

The frozen values are
\begin{equation}
    w_{\mathrm{1D}}^* = 0.72, \qquad w_{\mathrm{SE3D}}^* = 0.28, \qquad
    \tau^* = 0.47022256531284556 \approx 0.4702,
    \label{eq:frozen_values}
\end{equation}
where the full-precision threshold is rounded to $0.4702$ when reported in
the final decision rule. The final probability and class prediction are
therefore
\begin{equation}
    p_{\mathrm{GHXTox}} = 0.72\, p_{1D}^{\mathrm{cal}} + 0.28\,
    p_{\mathrm{SE3D}}, \qquad
    \hat{y} = \mathbb{I}\left[p_{\mathrm{GHXTox}} \geq 0.4702225653\right].
    \label{eq:final_prediction}
\end{equation}

A second-level leave-one-fold-out analysis is additionally performed as a
stability diagnostic. In each iteration, $w$ and $\tau$ are selected using
four OOF folds and applied to the remaining fold. This analysis assesses
whether the fusion parameters are disproportionately dependent on a
particular fold. However, because every upstream prediction stream is not
fully regenerated within a separately nested outer loop, this diagnostic is
not described as strict nested external validation.

The final values of 0.72/0.28 and 0.47022256531284556 are frozen from the complete
training OOF table. The corresponding fused OOF metrics describe
training-stage model selection rather than benchmark test-set performance.
After freezing, the same calibration parameters, fusion weights, and
threshold are applied unchanged to Test1 and Test2:
\begin{equation}
    \left(a^*, b^*, w^*, \tau^*\right)_{\mathrm{Test}} =
    \left(a^*, b^*, 0.72, 0.47022256531284556\right)_{\mathrm{Training\;OOF}}.
    \label{eq:test_protocol}
\end{equation}
Neither Test1 nor Test2 labels are used for recalibration or parameter
adjustment. The frozen parameters, input prediction files, fusion space, and
decision threshold are stored in machine-readable artifacts to permit
auditing of the final evaluation protocol.

\subsection{Evaluation Metrics}
To comprehensively evaluate GHXTox for peptide toxicity prediction, model
performance is assessed from three complementary perspectives: classification
performance at a fixed operating threshold, ranking ability based on
continuous probability scores, and probabilistic calibration. Reliance on a
single metric may conceal important differences in minority-class
recognition, false-positive control, or probability reliability. We therefore
jointly report confusion-matrix metrics, the area under the receiver
operating characteristic curve (AUROC), the area under the precision--recall
curve (AUPRC), the Brier score, and the expected calibration error (ECE).

Let the evaluation set contain $N$ peptides. The ground-truth label of peptide
$i$ is
\begin{equation}
    y_i \in \{0, 1\}, \qquad p_i = \mathbb{P}\left(y_i = 1 \mid
    \mathcal{S}_i\right),
    \label{eq:eval_setup}
\end{equation}
where $y_i = 1$ denotes a toxic peptide, $y_i = 0$ denotes a non-toxic
peptide, and $\mathcal{S}_i$ denotes the peptide sequence. The binary
prediction is obtained using a fixed decision threshold $\tau$:
\begin{equation}
    \hat{y}_i = \mathbb{I}\left(p_i \geq \tau\right).
    \label{eq:decision_threshold}
\end{equation}
For the final GHXTox model, the threshold is frozen using the complete
training OOF table, as described in Section 2.5.4. Its full-precision value
is $\tau^* = 0.47022256531284556$, which is rounded to $\tau = 0.4702$ for
presentation in the manuscript and result tables. The same frozen threshold
is applied unchanged to Test1 and Test2, without dataset-specific retuning.
Threshold-dependent metrics are calculated from $\hat{y}_i$, whereas AUROC,
AUPRC, the Brier score, and ECE are calculated directly from the continuous
probability $p_i$.

Because Test2 contains only 46 toxic peptides and 536 non-toxic peptides, its
class distribution is strongly imbalanced. Accuracy alone may therefore
provide an incomplete description of toxic-peptide recognition. Balanced
accuracy (BACC) and the Matthews correlation coefficient (MCC) are emphasized
as balanced summary measures; sensitivity (SN) and specificity (SP) are
reported to expose the trade-off between missed toxic peptides and false
alarms; and AUPRC and calibration metrics are reported together with AUROC.

\subsubsection{Threshold-Dependent Classification Metrics}
At a fixed decision threshold, the predictions can be summarized by four
confusion-matrix counts:
\begin{equation}
    \begin{aligned}
        \mathrm{TP} &= \sum_{i=1}^{N}\mathbb{I}\left(y_i=1,\;
        \hat{y}_i=1\right), &
        \mathrm{TN} &= \sum_{i=1}^{N}\mathbb{I}\left(y_i=0,\;
        \hat{y}_i=0\right),\\
        \mathrm{FP} &= \sum_{i=1}^{N}\mathbb{I}\left(y_i=0,\;
        \hat{y}_i=1\right), &
        \mathrm{FN} &= \sum_{i=1}^{N}\mathbb{I}\left(y_i=1,\;
        \hat{y}_i=0\right),
    \end{aligned}
    \label{eq:confusion}
\end{equation}
where TP denotes toxic peptides correctly identified as toxic, TN denotes
non-toxic peptides correctly identified as non-toxic, FP denotes non-toxic
peptides incorrectly predicted as toxic, and FN denotes toxic peptides
incorrectly predicted as non-toxic.

\textbf{Accuracy.}
Accuracy measures the proportion of correctly classified samples:
\begin{equation}
    \mathrm{ACC} = \frac{\mathrm{TP} + \mathrm{TN}}
    {\mathrm{TP} + \mathrm{TN} + \mathrm{FP} + \mathrm{FN}}.
    \label{eq:acc}
\end{equation}
ACC provides an overall correctness measure but may be dominated by the
majority class in strongly imbalanced datasets. A model biased toward
non-toxic predictions may still achieve high ACC when non-toxic peptides
greatly outnumber toxic peptides. ACC is therefore not used as the sole
measure of predictive performance.

\textbf{Sensitivity and specificity.}
Sensitivity, also referred to as recall or the true-positive rate, measures
the fraction of toxic peptides correctly identified:
\begin{equation}
    \mathrm{SN} = \frac{\mathrm{TP}}{\mathrm{TP} + \mathrm{FN}}, \qquad
    \mathrm{SP} = \frac{\mathrm{TN}}{\mathrm{TN} + \mathrm{FP}}.
    \label{eq:sn_sp}
\end{equation}
A high SN indicates that relatively few toxic peptides are missed. In a
toxicity-screening context, FN corresponds to potentially hazardous peptides
that are not detected, making SN an important measure of toxic-class
coverage. Specificity, or the true-negative rate, measures the fraction of
non-toxic peptides correctly rejected. A high SP indicates that relatively
few non-toxic candidates are falsely labeled as toxic. SN and SP should be
interpreted jointly: lowering the decision threshold generally increases SN
but may increase FP and reduce SP, whereas increasing the threshold may
produce the opposite trade-off.

\textbf{Precision and F1 score.}
Precision measures the proportion of predicted toxic peptides that are truly
toxic:
\begin{equation}
    \mathrm{Precision} = \frac{\mathrm{TP}}{\mathrm{TP} + \mathrm{FP}},
    \qquad
    F_1 = \frac{2 \cdot \mathrm{Precision} \cdot \mathrm{SN}}
    {\mathrm{Precision} + \mathrm{SN}} =
    \frac{2\,\mathrm{TP}}{2\,\mathrm{TP} + \mathrm{FP} + \mathrm{FN}}.
    \label{eq:precision_f1}
\end{equation}
Precision reflects the reliability of positive toxicity predictions. Because
Precision depends on the prevalence of toxic peptides, comparisons across
datasets with different class ratios should be interpreted together with the
corresponding sample composition. The F1 score balances positive-prediction
purity and toxic-peptide recall; however, it does not explicitly include TN
and therefore does not fully describe performance on the non-toxic class.

\textbf{Balanced accuracy.}
Balanced accuracy is the arithmetic mean of sensitivity and specificity:
\begin{equation}
    \mathrm{BACC} = \frac{\mathrm{SN} + \mathrm{SP}}{2}.
    \label{eq:bacc}
\end{equation}
BACC assigns equal importance to the toxic and non-toxic classes, preventing
the larger class from directly dominating the final score. A model performing
well on one class but poorly on the other will therefore have a noticeably
lower BACC than ordinary ACC.

\textbf{Matthews correlation coefficient.}
The Matthews correlation coefficient incorporates all four elements of the
confusion matrix:
\begin{equation}
    \mathrm{MCC} = \frac{\mathrm{TP} \cdot \mathrm{TN} -
    \mathrm{FP} \cdot \mathrm{FN}}
    {\sqrt{(\mathrm{TP}+\mathrm{FP})(\mathrm{TP}+\mathrm{FN})
    (\mathrm{TN}+\mathrm{FP})(\mathrm{TN}+\mathrm{FN})}},
    \label{eq:mcc}
\end{equation}
with $-1 \leq \mathrm{MCC} \leq 1$. An MCC of 1 indicates perfect agreement,
0 indicates no overall correlation between predictions and labels, and $-1$
indicates complete disagreement. Because MCC jointly accounts for TP, TN, FP,
and FN, it was used as the primary summary metric for fixed-threshold
classification and as the primary objective for OOF fusion-parameter
selection \cite{matthews1975mcc}.

Numerical safeguards were applied when metric denominators were zero. All
threshold-dependent metrics were calculated using the same OOF-frozen GHXTox
decision threshold, consistent with the freezing protocol described in
Section 2.5.4; no metric-specific or test-set-specific threshold optimization
was performed.

\subsubsection{Ranking and Probability-Calibration Metrics}
Fixed-threshold metrics describe classification at one operating point but do
not fully characterize sample ranking across all possible thresholds or
determine whether the predicted probabilities have reliable numerical
meaning. We therefore additionally report AUROC, AUPRC, the Brier score, and
ECE-10. AUROC and AUPRC evaluate probability ranking, whereas the Brier score
and ECE evaluate probabilistic error and calibration.

\textbf{ROC curve and AUROC.}
For a varying probability threshold $t$, the true-positive rate and the
false-positive rate are
\begin{equation}
    \mathrm{TPR}(t) = \frac{\mathrm{TP}(t)}{\mathrm{TP}(t) +
    \mathrm{FN}(t)}, \qquad
    \mathrm{FPR}(t) = \frac{\mathrm{FP}(t)}{\mathrm{FP}(t) +
    \mathrm{TN}(t)}.
    \label{eq:tpr_fpr}
\end{equation}
The ROC curve plots $\mathrm{TPR}(t)$ against $\mathrm{FPR}(t)$ over all
thresholds, and its area is
\begin{equation}
    \mathrm{AUROC} = \int_{0}^{1} \mathrm{TPR}(u)\, du,
    \label{eq:auroc}
\end{equation}
where $u$ denotes the false-positive rate. AUROC measures the ability of the
model to rank a randomly selected toxic peptide above a randomly selected
non-toxic peptide. Values approaching one indicate strong ranking
performance, whereas values near 0.5 generally indicate approximately random
ranking. AUROC is independent of the fixed decision threshold and therefore
complements the threshold-dependent metrics. Although AUROC is insensitive to
class prevalence, it may underemphasize positive-class retrieval quality and
the practical false-positive burden under severe class imbalance. AUPRC is
therefore reported as a complementary ranking metric. Recent studies in
peptide toxicity prediction commonly report ROC and PR analyses together
\cite{zhang2025toxmsrc,liang2026toxpltc,guan2025toxipep,jiao2025structoxnet}.

\textbf{Precision--Recall curve and AUPRC.}
For threshold $t$, Precision and Recall are defined as
\begin{equation}
    \mathrm{Precision}(t) = \frac{\mathrm{TP}(t)}{\mathrm{TP}(t) +
    \mathrm{FP}(t)}, \qquad
    \mathrm{Recall}(t) = \frac{\mathrm{TP}(t)}{\mathrm{TP}(t) +
    \mathrm{FN}(t)}.
    \label{eq:precision_recall}
\end{equation}
The PR curve describes the trade-off between positive-prediction purity and
toxic-peptide recall as the threshold varies. In the current implementation,
AUPRC is calculated as average precision through
\texttt{average\_precision\_score}. Let $P_n$ and $R_n$ denote Precision and
Recall at the $n$-th score threshold after samples have been ranked in
descending order; the reported value is
\begin{equation}
    \mathrm{AUPRC}_{\mathrm{reported}} \coloneqq \mathrm{AP} = \sum_{n}
    \left(R_n - R_{n-1}\right) P_n.
    \label{eq:average_precision}
\end{equation}
Throughout this study, the metric is reported as AUPRC, while its numerical
value is computed using average precision. This definition weights Precision
by the increase in Recall and is not necessarily identical to a trapezoidal
numerical integral of the PR curve. AUPRC focuses on retrieval of the
positive toxic class and is particularly informative when toxic peptides are
relatively rare. Because the chance-level PR baseline is determined by the
positive-class prevalence, AUPRC values for Test1 and Test2 should be
interpreted within their respective class distributions rather than treated
as directly comparable measures of dataset difficulty. AUROC and AUPRC are
reported only when both classes are present in the evaluation subset;
otherwise, the implementation returns an undefined value.

\textbf{Brier score.}
The Brier score is the mean squared difference between the predicted
probability and the binary ground-truth label \cite{brier1950}:
\begin{equation}
    \mathrm{Brier} = \frac{1}{N} \sum_{i=1}^{N}
    \left(p_i - y_i\right)^2.
    \label{eq:brier}
\end{equation}
Its range is $0 \leq \mathrm{Brier} \leq 1$, with lower values indicating
better probabilistic predictions. The Brier score penalizes both incorrect
predictions and excessive confidence: assigning a probability of 0.01 to a
truly toxic peptide incurs a larger Brier error than assigning a probability
of 0.40. The Brier score reflects both probability calibration and
discriminative performance and should not be interpreted as a pure
calibration measure; it is therefore considered jointly with AUROC, AUPRC,
and ECE.

\textbf{Ten-bin expected calibration error.}
To further quantify the agreement between predicted confidence and observed
toxicity frequency, we calculate the expected calibration error using ten
equal-width probability bins, denoted ECE-10 \cite{guo2017calibration}.
The interval $[0,1]$ is divided into $M = 10$ bins, and for each
non-empty bin the mean predicted confidence and the observed toxic-peptide
frequency are
\begin{equation}
    \begin{aligned}
        \mathcal{B}_m &= \left\{i : p_i \in \left[\frac{m-1}{M},\;
        \frac{m}{M}\right)\right\}, \quad m = 1, \ldots, M,\\
        \mathrm{conf}\left(\mathcal{B}_m\right) &=
        \frac{1}{|\mathcal{B}_m|} \sum_{i \in \mathcal{B}_m} p_i,\\
        \mathrm{freq}\left(\mathcal{B}_m\right) &=
        \frac{1}{|\mathcal{B}_m|} \sum_{i \in \mathcal{B}_m} y_i,
    \end{aligned}
    \label{eq:ece_bins}
\end{equation}
with the final bin including $p_i = 1$. ECE-10 is defined as
\begin{equation}
    \mathrm{ECE}_{10} = \sum_{m=1}^{10} \frac{|\mathcal{B}_m|}{N}\;
    \left|\mathrm{conf}\left(\mathcal{B}_m\right) -
    \mathrm{freq}\left(\mathcal{B}_m\right)\right|.
    \label{eq:ece}
\end{equation}
Values closer to zero indicate better agreement between predicted confidence
and observed toxicity frequency; for example, a well-calibrated group of
predictions near 0.8 should contain approximately 80\% toxic peptides.
Because ECE depends on the binning scheme, the number and boundaries of the
bins remain fixed across models and datasets.

\textbf{Joint interpretation of the metrics.}
No individual metric is sufficient to characterize peptide-toxicity
prediction; the reported measures are therefore interpreted jointly. ACC
summarizes overall correctness; SN and SP quantify toxic-peptide misses and
non-toxic false positives, respectively; BACC evaluates class-balanced
recognition; MCC provides a comprehensive correlation-based measure using all
four confusion-matrix terms; AUROC and AUPRC assess ranking performance; and
the Brier score and ECE-10 characterize probabilistic reliability. Exact
numerical tables are used as the primary form of model comparison, with ROC
and PR curves presented together. Radar charts, when included, are treated
only as secondary summaries and do not replace exact metric tables or
uncertainty analysis.

\subsection{Robustness, Uncertainty, and Statistical Analysis}
A single performance point estimate cannot fully characterize the reliability
of a deep learning predictor. Model performance may be affected by random
initialization, finite evaluation samples, class imbalance, predicted-
structure quality, and distribution shift. The GHXTox evaluation framework
therefore prespecifies five complementary perspectives: multi-seed
reproducibility, class-stratified bootstrap confidence intervals, paired
significance testing between frozen models, training-OOF conformer
information-ceiling diagnostics, and subgroup and structure-quality analysis.
For the retained
multi-conformer benchmark output, inferential statistics are reported only
after they have been regenerated from the corresponding sample-level
prediction vectors; intervals from preceding model variants are not reused.

All analyses in this section are conducted after the model architecture,
checkpoints, probability-calibration parameters, fusion weights, and decision
threshold have been frozen. No ProtT5, Mamba, EGNN, chemical-site, or
conformer-residual parameters are updated during statistical evaluation, and
the threshold is not modified according to the resulting confidence intervals
or significance tests. This design explicitly separates model development
from frozen-model evaluation.

Three sources of uncertainty are distinguished. Multi-seed analysis
characterizes variation arising from stochastic optimization; bootstrap
intervals quantify finite-sample uncertainty conditional on a frozen
predictor; and subgroup analyses examine whether predictive performance is
concentrated within particular sequence-length, physicochemical, or
structure-confidence ranges. These analyses provide complementary evidence
and should not be interpreted as interchangeable measures of uncertainty.

\subsubsection{Multi-Seed Reproducibility and Prediction Dispersion}
Neural-network training may vary because of parameter initialization,
mini-batch ordering, dropout masks, and mixed-precision computation. To
assess the sensitivity of the 1D expert to these factors, three LoRA-ProtT5
and bidirectional Mamba models are trained using the random seeds
$\mathcal{R} = \{42, 123, 2025\}$, while keeping the data partitions,
architecture, training duration, and all hyperparameters unchanged. The
current 1D expert uses an equal-weight probability average of the three
full-data models rather than selecting the best-performing seed post hoc. The
top-level model retains a single frozen 3D structural expert: the evaluated
multi-checkpoint 3D ensemble did not satisfy the prespecified promotion
criterion and was therefore not adopted in the current default GHXTox model.

For an evaluation metric $M$, let $M^{(s)}$ denote the result obtained using
seed $s$. The multi-seed mean and the sample standard deviation are
\begin{equation}
    \begin{aligned}
        \bar{M} &= \frac{1}{S} \sum_{s=1}^{S} M^{(s)}, &
        \mathrm{SD}(M) &= \sqrt{\frac{1}{S-1} \sum_{s=1}^{S}
        \left(M^{(s)} - \bar{M}\right)^2}, \qquad S = 3,
    \end{aligned}
    \label{eq:seed_metrics}
\end{equation}
and seed-level reproducibility may therefore be reported as
$\bar{M} \pm \mathrm{SD}(M)$. All seed-specific models must be evaluated
using the same frozen threshold, or using thresholds selected without access
to the corresponding test set; a test-specific optimal threshold must not be
selected separately for each seed.

Prediction dispersion can also be quantified at the sample level. For peptide
$i$, let $p_{i,\mathrm{1D}}^{(s)}$ denote the uncalibrated probability from
seed $s$. The seed-averaged probability and the seed dispersion are
\begin{equation}
    \begin{aligned}
        \bar{p}_{i,\mathrm{1D}}^{\mathrm{raw}} &= \frac{1}{S}
        \sum_{s=1}^{S} p_{i,\mathrm{1D}}^{(s)}, &
        u_{i,\mathrm{seed}} &= \sqrt{\frac{1}{S-1} \sum_{s=1}^{S}
        \left(p_{i,\mathrm{1D}}^{(s)} -
        \bar{p}_{i,\mathrm{1D}}^{\mathrm{raw}}\right)^2},
    \end{aligned}
    \label{eq:seed_dispersion}
\end{equation}
where $\bar{p}_{i,\mathrm{1D}}^{\mathrm{raw}}$ denotes the uncalibrated
three-seed average used at inference, which is distinguished from the
training-time OOF average $\bar{p}_{i,\mathrm{1D}}^{\mathrm{OOF}}$ defined in
Section 2.5.1. A large $u_{i,\mathrm{seed}}$ indicates that the prediction is
sensitive to optimization randomness and may be used to flag boundary cases
or expert-disagreement samples. It is, however, an empirical measure of
training variability rather than a Bayesian posterior uncertainty estimate or
a direct probability of prediction error.

The final 1D pathway follows the order defined in Section 2.3.4:
\begin{equation*}
    \left\{p_{i,\mathrm{1D}}^{(s)}\right\}_{s=1}^{3}
    \longrightarrow \bar{p}_{i,\mathrm{1D}}^{\mathrm{raw}}
    \longrightarrow \text{OOF-fitted Platt calibration}
    \longrightarrow p_{i,\mathrm{1D}}^{\mathrm{cal}}.
\end{equation*}
Seed-specific metrics characterize training stability, whereas the calibrated
three-seed average is used for top-level 1D/3D fusion. The best individual
seed is not substituted for the expected multi-seed result. The final
calibrated ensemble prediction is a deterministic function of the three
trained seed models; seed-level variability is summarized by
$u_{i,\mathrm{seed}}$ rather than being treated as zero. Accordingly,
variability is reported across the constituent seed models rather than by
assigning a standard deviation of zero to their averaged prediction.

\subsubsection{Class-Stratified Bootstrap Confidence Intervals}
Performance metrics calculated on a finite test set are themselves estimates.
This issue is particularly important for Test2, where a small number of
changes in TP or FN can substantially alter SN, BACC, and MCC. We therefore
use class-stratified percentile bootstrap resampling to calculate 95\%
confidence intervals \cite{efron1993bootstrap}.

Let the positive- and negative-class index sets be
\begin{equation}
    \begin{aligned}
        \mathcal{I}_{+} &= \{i : y_i = 1\}, &
        \mathcal{I}_{-} &= \{i : y_i = 0\},\\
        \mathbf i_{+}^{(b)} &\sim \mathrm{SampleWithReplacement}\left(
        \mathcal{I}_{+},\; |\mathcal{I}_{+}|\right), &
        \mathbf i_{-}^{(b)} &\sim \mathrm{SampleWithReplacement}\left(
        \mathcal{I}_{-},\; |\mathcal{I}_{-}|\right),\\
        \mathbf i^{(b)} &= \mathrm{Concat}\left(\mathbf i_{+}^{(b)},\;
        \mathbf i_{-}^{(b)}\right),
    \end{aligned}
    \label{eq:bootstrap_resample}
\end{equation}
where the bootstrap index vectors preserve repeated indices. This procedure
preserves the original number of positive and negative samples in every
replicate, avoiding extreme class-ratio fluctuations or single-class
replicates when the positive class is small. The validation implementation
independently resamples the positive and negative index sets with
replacement.

For metric $M$, the bootstrap estimate is
\begin{equation}
    M^{(b)} = M\left(\mathbf y_{\mathbf i^{(b)}},\;
    \mathbf p_{\mathbf i^{(b)}}\right), \qquad b = 1, \ldots, B.
    \label{eq:bootstrap_metric}
\end{equation}
The primary test-set analysis used $B = 5000$ bootstrap replicates, whereas
subgroup-difference analyses used 1000 replicates. The percentile confidence
interval is
\begin{equation}
    \mathrm{CI}_{95\%}(M) = \left[Q_{0.025}\left(\left\{M^{(b)}\right\}_{b=1}^{B}
    \right),\; Q_{0.975}\left(\left\{M^{(b)}\right\}_{b=1}^{B}\right)\right],
    \label{eq:bootstrap_ci}
\end{equation}
where $Q_{\alpha}$ is the empirical $\alpha$-quantile.

For ACC, BACC, SN, SP, F1, and MCC, each bootstrap replicate uses the
original frozen model threshold; the threshold is not reoptimized within any
resampled data set. AUROC, AUPRC, and the Brier score are calculated directly
from the resampled continuous probabilities.

These intervals quantify finite-sample uncertainty conditional on the frozen
predictor and the observed evaluation distribution; they do not incorporate
training-seed variability or distribution shift. Bootstrap uncertainty is
therefore interpreted separately from the multi-seed analysis in
Section 2.7.1 and the subgroup analysis in Section 2.7.5.

\subsubsection{Paired Bootstrap and Randomization Significance Testing}
When the current GHXTox model is compared with the previous default or
another frozen internal baseline, both models produce predictions for the
same peptides; the sample-level predictions are therefore paired. Comparing
two independent confidence intervals would ignore this dependence. We instead
analyze the metric difference directly on aligned samples.

Let $M_A$ and $M_B$ denote the metric values of frozen models $A$ and $B$,
respectively, where $B$ is the current model and $A$ is the frozen baseline.
The observed difference is
\begin{equation}
    \Delta_M^{\mathrm{obs}} = M_B - M_A.
    \label{eq:obs_difference}
\end{equation}
For BACC, SN, SP, F1, MCC, AUROC, and AUPRC, a positive difference favors
model $B$; for lower-is-better measures such as the Brier score and ECE, a
negative difference indicates improvement.

\textbf{Paired class-stratified bootstrap.}
For every bootstrap replicate, the same class-stratified index vector
$\mathbf i^{(b)}$ defined in Section 2.7.2 is applied to both models:
\begin{equation}
    \Delta_M^{(b)} = M\left(\mathbf y_{\mathbf i^{(b)}},\;
    \mathbf p_{B,\mathbf i^{(b)}}\right) - M\left(\mathbf y_{\mathbf i^{(b)}},\;
    \mathbf p_{A,\mathbf i^{(b)}}\right),
    \label{eq:paired_bootstrap}
\end{equation}
where using identical indices preserves the paired sample structure. The 95\%
percentile interval is
\begin{equation}
    \mathrm{CI}_{95}\left(\Delta_M\right) = \left[Q_{0.025}\left(\left\{
    \Delta_M^{(b)}\right\}_{b=1}^{B}\right),\; Q_{0.975}\left(\left\{
    \Delta_M^{(b)}\right\}_{b=1}^{B}\right)\right].
    \label{eq:paired_ci}
\end{equation}
For a higher-is-better metric, an interval entirely above zero supports a
consistent positive difference; for a lower-is-better metric such as the
Brier score, improvement corresponds to an interval entirely below zero.
Using a metric-direction coefficient $s_M$, where $s_M = +1$ for
higher-is-better metrics and $s_M = -1$ for lower-is-better metrics, the
bootstrap probability of improvement is
\begin{equation}
    P_{\mathrm{boot}}^{\mathrm{improve}} = \frac{1}{B}
    \sum_{b=1}^{B} \mathbb{I}\left(s_M\, \Delta_M^{(b)} > 0\right).
    \label{eq:bootstrap_improve}
\end{equation}
This quantity is the proportion of bootstrap replicates in which the current
model outperforms the baseline; it is not a classical hypothesis-test
$p$-value and should not be reported using an undifferentiated $P$ label.

\textbf{Paired randomization test.}
A within-sample exchange randomization test is additionally used to evaluate
the null hypothesis that the two frozen model outputs are exchangeable for
each peptide \cite{ojala2010permutation}. For randomization replicate
$r$, $z_i^{(r)} \sim \mathrm{Bernoulli}(0.5)$ is independently drawn for each
sample; if $z_i^{(r)} = 1$, the probabilities and frozen-threshold decisions
of models $A$ and $B$ are exchanged for sample $i$, and otherwise they remain
unchanged. The resulting null difference is
\begin{equation}
    \Delta_{M,0}^{(r)} = M_B^{(r)} - M_A^{(r)}.
    \label{eq:null_difference}
\end{equation}
A two-sided paired randomization test was performed by comparing the
absolute randomized metric differences with the absolute observed difference:
\begin{equation}
    p_{\mathrm{rand}} = \frac{1 + \sum_{r=1}^{R}
    \mathbb{I}\left(\left|\Delta_{M,0}^{(r)}\right| \geq
    \left|\Delta_M^{\mathrm{obs}}\right|\right)}{R + 1}.
    \label{eq:rand_pvalue}
\end{equation}
The primary paired analysis uses $B = 5000$ and $R = 5000$ replicates with a
fixed seed of 2026. The implementation exchanges both continuous scores and
the corresponding frozen-threshold decisions within each sample.

MCC and BACC are treated as prespecified primary summary measures, whereas
SN, SP, AUROC, AUPRC, and the Brier score provide complementary evidence.
Because several correlated metrics are examined for the same model
comparison, secondary $p$-values are interpreted descriptively rather than as
multiple independent discoveries.

Finally, Test1 and Test2 have been repeatedly observed during the historical
development of the project. When recomputed for the retained prediction
vectors, paired confidence intervals and randomization $p$-values on these
sets provide retrospective evidence for differences between frozen models,
rather than fully independent prospective confirmation; this limitation is
explicitly recorded by the project validation implementation.

\subsubsection{Training-OOF Conformer Information-Ceiling Diagnostics}
To distinguish limited conformer diversity from limited use of that diversity
by the frozen predictor, we conducted two post hoc diagnostics using only the
group-aware training OOF folds. Neither diagnostic accessed Test1 or Test2,
altered the retained full2795 checkpoint, or participated in model selection.

First, each of the eight conformers available for 2{,}795 covered training
peptides was evaluated with the corresponding frozen fold-specific structural
expert, yielding 22{,}360 held-view predictions. For each peptide, we summarized
the maximum pairwise conformer RMSD, mean residue RMSF, and range of the eight
structural-expert probabilities. Label-free fixed aggregators, including the
mean, median, trimmed mean, confidence- and centrality-weighted means, and a
structural-medoid rule, were compared with the frozen single-structure output
under the unchanged top-level fusion and decision threshold. A label-aware
choice among the single-structure and per-conformer outputs was computed only as a nondeployable
upper-bound diagnostic and was not treated as a model result.

Second, pooled spatial embeddings were extracted for every conformer using the
same fold-specific frozen expert. A compact 50{,}113-parameter
permutation-invariant residual head
was fitted under fold-local train/held separation, while all upstream encoders,
the 1D probability stream, the top-level fusion weights, and the decision
threshold remained fixed. The five held-fold prediction vectors were
concatenated once for paired OOF evaluation and class-stratified bootstrap
analysis. Read-only diagnostics also quantified the Pearson correlation between
the 1D and base 3D probabilities, the correlation between the learned residual
and the correction required by the observed label, and the fraction of residual
directions that moved the fused probability toward that label. These quantities
were used to characterize a possible representation--label alignment bottleneck,
not to tune the candidate head.

\subsubsection{Subgroup and Structure-Quality Analysis}
Aggregate test metrics may conceal performance heterogeneity across peptide
lengths, physicochemical properties, and structure-quality ranges. This issue
is especially important for structure-dependent models: an overall gain may
be concentrated among high-pLDDT peptides and may not generalize to short
peptides, uncertain structures, or samples with incomplete chemical-site
coverage. We therefore conduct prespecified subgroup analyses without
modifying the frozen model or threshold.

The primary current-model length analysis uses 11--15, 16--30, and 31--50
residues. General structural diagnostics additionally stratify samples by
mean pLDDT (<0.55, 0.55--0.70, and $\geq$0.70), net charge (negative, neutral,
and positive), mean hydropathy (<0, 0--1, and $\geq$1), cysteine count (0,
1--2, and $\geq$3), and chemical-site coverage (<0.80, 0.80--0.95, and
$\geq$0.95). Chemical-site coverage is defined as the fraction of residues
containing at least one valid chemical site:
\begin{equation}
    r_{\mathrm{site}} = \frac{1}{L} \sum_{i=1}^{L}
    \mathbb{I}\left(\text{residue } i \text{ contains at least one valid
    chemical site}\right).
    \label{eq:site_coverage}
\end{equation}
The validation implementation also records the minimum pLDDT, the low-pLDDT
residue fraction, the net charge, the mean hydropathy, and the cysteine
count.

To reduce unstable conclusions from extremely small subsets, a stratum is
reported only when $N_{\mathrm{stratum}} \geq 20$ and
$N_{+,\mathrm{stratum}} \geq 5$ and $N_{-,\mathrm{stratum}} \geq 5$. Strata
that failed these reporting criteria were retained in the audit output with
their sample counts and the reason for omission, but were excluded from
metric-difference interpretation. These requirements match the current
subgroup-validation implementation.

For each eligible stratum, BACC, SN, SP, F1, and MCC are calculated using the
same threshold frozen on the full training OOF predictions; AUROC, AUPRC, and
the Brier score are computed from the original continuous probabilities. No
stratum-specific threshold optimization is performed.

When comparing the current model with a frozen baseline, the subgroup
difference is
\begin{equation}
    \Delta M_g = M_{B,g} - M_{A,g},
    \label{eq:subgroup_delta}
\end{equation}
where $g$ denotes a stratum. Class-stratified bootstrap intervals are
calculated for the MCC and AUPRC differences, using 1000 subgroup resamples
in the current validation protocol.

Subgroup results are treated as heterogeneity diagnostics rather than new
model-selection evidence. Because multiple variables and bins are examined,
isolated positive or negative findings are interpreted together with the
sample size, the positive-class count, the neighboring-stratum trends, and
the confidence intervals.

In particular, pLDDT is a confidence estimate produced by a
structure-prediction model. Better performance in a high-pLDDT subgroup does
not establish correspondence to experimentally correct structures, and poorer
performance at low pLDDT does not imply that the corresponding residues lack
biological relevance. GHXTox uses continuous confidence weighting rather than
excluding low-confidence samples; subgroup analysis is used to characterize
applicability boundaries rather than to justify hard filtering.

\subsection{Interpretability Analysis Protocol}
Interpretability analyses are defined on the frozen model and do not
participate in feature selection, threshold selection, or fusion-weight
optimization. Unless otherwise stated, probability changes are computed with
all non-target inputs held fixed, summarized separately for toxic and
non-toxic peptides, and accompanied by peptide-level bootstrap intervals.
The analyses therefore describe model sensitivity and expert complementarity;
they are not interpreted as evidence of a causal toxicophore.

\subsubsection{Fixed-Representation Residue Occlusion and Seed Stability}
Residue-level sensitivity was evaluated for a prespecified 16-peptide cohort
selected independently of the attribution values. Within each test set and
each expert-outcome group (both experts correct, only the 1D expert correct,
only the 3D expert correct, or both experts incorrect), one toxic and one
non-toxic peptide nearest the frozen fusion threshold were selected. Ties were
resolved by source index. This rule prevents favorable cases from being chosen
after inspecting the attribution maps.

For each of the three frozen 1D checkpoints, task-adapted ProtT5 hidden states
were computed in the original 16-record inference batches, thereby preserving
the deployment-time padding configuration and the boundary context of the
backward convolution.
One residue-aligned, post-contextual ProtT5 vector was then set to zero while
all other contextual vectors, the sequence length, bidirectional Mamba layers,
pooling operation, and classifier parameters were held fixed. The per-seed
occlusion effect was defined as the unoccluded logit minus the occluded logit.
For every perturbed state, the three raw seed probabilities were averaged
before applying the frozen OOF-fitted Platt transform. The displayed signed
attribution was the calibrated probability change normalized by the
within-peptide $L_1$ norm. Stability was summarized by residue-wise sign
agreement and the pairwise Spearman correlation of absolute logit effects
across seeds. This fixed-representation experiment quantifies the sensitivity
of the downstream 1D layers to contextual residue vectors; it is neither a
token-mutation experiment nor an Integrated Gradients analysis, and it does
not test residue-level biological causality.

\subsubsection{Chemical-Site Deletion Analysis}
For the same prespecified cohort, each valid modeled chemical site was ablated
by setting its availability mask, type vector, coordinate, orientation, and
orientation-validity entry to zero. All non-target sequence features, residue
geometry, pLDDT values, other chemical sites, conformer descriptors, model
weights, and operating thresholds were held fixed. The complete frozen 3D
forward pass was recomputed, including the conformer residual when an audited
cache entry was available. Complementary perturbations deleted all sites
carrying one chemical category or all site slots assigned to one residue. We
report $\Delta p_{3D}=p_{3D}^{\mathrm{base}}-p_{3D}^{\mathrm{deleted}}$ and,
with the 1D probability fixed, the induced fusion change
$\Delta p_{\mathrm{fused}}=0.28\Delta p_{3D}$. Because the modeled sites and
candidate spatial relations are derived from predicted structures, these
quantities describe frozen-model sensitivity rather than confirmed chemical
bonds or intervention effects.

\subsubsection{pLDDT--Gate Association Analysis}
The unperturbed structural expert was additionally run on all 1{,}126 Test1
and 582 Test2 peptides. For each peptide, we recorded the mean pLDDT, learned
global gate, and the within-peptide Spearman correlation between residue-level
pLDDT and node gates. Associations were summarized separately by test set and
for the pooled descriptive sample. Ninety-five percent bootstrap percentile
intervals were obtained from 2{,}000 peptide-level resamples with seed 2026.
This analysis characterizes the association between predicted-structure
confidence and learned gate values; it does not establish structural accuracy,
feature attribution, or biological causality.

\subsubsection{Expert Agreement and Conflict Analysis}
Each benchmark peptide is assigned to one of four predeclared groups:
both experts correct, only the 1D expert correct, only the 3D expert correct,
or both experts incorrect. We compare the sequence probability, structural
probability, fused probability, peptide length, class label, and structural
confidence across these groups. Expert outcomes use the training-derived
standalone thresholds of 0.392 for the 1D expert and 0.677819 for the 3D
expert; the fused decision uses the separately frozen threshold of
0.47022256531284556. Cases rescued by structural evidence are inspected
alongside cases in which structural evidence reduces an otherwise correct 1D
score. This analysis is performed after freezing all model parameters and is
used to describe complementarity rather than to retune the fusion rule on
Test1 or Test2.

% ============================================================================
%  3 RESULTS AND DISCUSSION
% ============================================================================
\clearpage
\section{RESULTS AND DISCUSSION}
\subsection{Overall Predictive Performance on the Benchmark Test Sets}
Using the OOF-frozen weight and threshold without dataset-specific retuning,
GHXTox correctly classified 1{,}068 of 1{,}126 Test1 peptides
(TP/TN/FP/FN = 287/781/25/33). This corresponds to ACC 0.9485, BACC 0.9329,
SN 0.8969, SP 0.9690, F1 0.9082, MCC 0.8726, AUROC 0.9811, and AUPRC
0.9675. On the strongly imbalanced Test2 set, GHXTox correctly classified 552
of 582 peptides (35/517/19/11), yielding ACC 0.9485, BACC 0.8627, SN 0.7609,
SP 0.9646, F1 0.7000, MCC 0.6746, AUROC 0.9433, and AUPRC 0.8043. The lower
Test2 BACC and F1, despite high ACC, reflect the difficulty of recovering the
46 toxic peptides and demonstrate why class-balanced and precision--recall
metrics are essential for this benchmark.

\subsubsection{Comparison with Existing Peptide-Toxicity Predictors}
Table~\ref{tab:overall_comparison} compares models reported on the same
benchmark partitions. Values for ToxMSRC and ToxPLTC are transcribed from the
corresponding publications or their accompanying result tables
\cite{zhang2025toxmsrc,liang2026toxpltc}; NR denotes an
unreported metric. GHXTox obtains the highest Test1 ACC, BACC, SP, MCC, AUROC,
and AUPRC among the compared rows, while ToxPLTC has higher Test1 SN. On
Test2, GHXTox obtains the highest ACC, SP, MCC, and reported AUPRC, but its
BACC and SN remain below both reference models. The supported conclusion is
therefore a stronger overall balance of specificity, correlation, and
positive-class ranking, not comprehensive superiority on every metric. Because
Test1 and Test2 had been observed historically, these differences remain
retrospective benchmark evidence and require confirmation on a new prefrozen
external set.

% ---- Table 3: overall comparison --------------------------------------------
\begin{table}[htbp]
    \centering
    \caption{Performance comparison on the benchmark Test1 and Test2 sets.}
    \label{tab:overall_comparison}
    \resizebox{\linewidth}{!}{%
    \begin{tabular}{llccccccc}
        \toprule
        Dataset & Model & ACC & BACC & SN & SP & MCC & AUROC & AUPRC \\
        \midrule
        \multirow{3}{*}{Test1}
        & ToxMSRC & 0.9405 & 0.9217 & 0.8781 & 0.9653 & 0.8520 & 0.9650 & NR \\
        & ToxPLTC & 0.9378 & 0.9302 & \textbf{0.9125} & 0.9479 & 0.8496 & 0.9709 & 0.9426 \\
        & GHXTox (current) & \textbf{0.9485} & \textbf{0.9329} & 0.8969 & \textbf{0.9690} & \textbf{0.8726} & \textbf{0.9811} & \textbf{0.9675} \\
        \midrule
        \multirow{3}{*}{Test2}
        & ToxMSRC & 0.9416 & 0.8689 & 0.7826 & 0.9552 & 0.6550 & 0.9430 & NR \\
        & ToxPLTC & 0.9261 & \textbf{0.8804} & \textbf{0.8261} & 0.9347 & 0.6197 & 0.9295 & 0.7353 \\
        & GHXTox (current) & \textbf{0.9485} & 0.8627 & 0.7609 & \textbf{0.9646} & \textbf{0.6746} & \textbf{0.9433} & \textbf{0.8043} \\
        \bottomrule
    \end{tabular}}
\end{table}

\subsubsection{Discrimination, Calibration, and Decision Characteristics}
The final fused probabilities have Brier scores of 0.0407 and 0.0430 on Test1
and Test2, respectively. These values summarize mean squared probability
error under the frozen inference protocol. ECE-10 is not reported here because
a consistently regenerated value is not available for both current
multi-conformer prediction files; retaining ECE values from the preceding
zero-conformer artifact would mix non-identical model outputs. Test2 AUPRC
(0.8043) is particularly informative: unlike AUROC, it
directly reflects the precision--recall trade-off for the rare toxic class.
The 11 false negatives on Test2 also delimit the current operating point; the
model's high specificity should not be interpreted as complete minority-class
coverage.

% ---- Figure 3: discrimination and calibration (placeholder) -----------------
\begin{figure}[htbp]
    \centering
    % TODO: \includegraphics[width=\linewidth]{figures/fig3_discrimination_calibration.pdf}
    \caption{Overall discrimination and calibration performance of GHXTox. (A, B) Test1 ROC and PR curves; (C, D) Test2 ROC and PR curves; (E, F) Test1 and Test2 reliability diagrams.}
    \label{fig:discrimination_calibration}
\end{figure}

\subsection{Complementarity of the 1D and 3D Experts}
The two experts are fused only after independent probability generation. The
1D branch provides the stronger standalone ranking signal, whereas the 3D
branch contributes confidence-weighted geometry and local chemical context.
This separation makes the gain attributable to complementary probability
evidence rather than duplicated use of an earlier fused prediction.

\subsubsection{Performance of the Individual Experts}
Table~\ref{tab:expert_comparison} reports threshold-free discrimination and
probability error for probability streams that are aligned with the retained
current configuration. Platt scaling preserves the 1D ranking but reduces its
Brier score. Relative to calibrated 1D alone, final fusion improves Test1
AUROC/AUPRC from 0.9799/0.9644 to 0.9811/0.9675 and reduces Brier from 0.0429
to 0.0407. On Test2, the corresponding values change from
0.9409/0.7831/0.0451 to 0.9433/0.8043/0.0430. The current fused predictor
therefore adds complementary structural ranking information while retaining a
lower probability error than the calibrated 1D stream. Historical
single-structure SE3D and equal-weight rows are not relabeled as current
component results because they were generated before activation of the audited
benchmark conformer cache.

% ---- Table 4: expert comparison ---------------------------------------------
\begin{table}[htbp]
    \centering
    \caption{Performance of the 1D probability streams and the retained current probability-level fusion. The final row for each benchmark uses the audited multi-conformer structural stream.}
    \label{tab:expert_comparison}
    \begin{tabular}{llccc}
        \toprule
        Dataset & Probability stream & AUROC & AUPRC & Brier \\
        \midrule
        \multirow{3}{*}{Test1}
        & 1D, uncalibrated & 0.9799 & 0.9644 & 0.0445 \\
        & 1D, OOF Platt calibrated & 0.9799 & 0.9644 & 0.0429 \\
        & Current fusion (0.72/0.28) & \textbf{0.9811} & \textbf{0.9675} & \textbf{0.0407} \\
        \midrule
        \multirow{3}{*}{Test2}
        & 1D, uncalibrated & 0.9409 & 0.7831 & 0.0491 \\
        & 1D, OOF Platt calibrated & 0.9409 & 0.7831 & 0.0451 \\
        & Current fusion (0.72/0.28) & \textbf{0.9433} & \textbf{0.8043} & \textbf{0.0430} \\
        \bottomrule
    \end{tabular}
\end{table}

\subsubsection{OOF-Frozen Probability Fusion}
The 0.72 sequence weight is an empirical result of the training-OOF search,
not a manually imposed assertion that sequence information is intrinsically
more important. Both expert weights and the full-precision decision threshold
of 0.47022256531284556 were fixed before Test1 or Test2 evaluation. The
held-out benchmark results therefore assess one prespecified operating point;
no benchmark-specific reweighting, threshold adjustment, or post hoc selection
between equal and unequal fusion is performed.

\subsubsection{Expert Agreement, Disagreement, and Error Correction}
Expert disagreement should be interpreted at the probability level rather than
as a causal statement about individual residues. The final weighting allows
structural evidence to alter borderline rankings while preventing the weaker
standalone 3D branch from dominating the decision. The modest observed effect
sizes and conditional structural coverage indicate that these corrections are
sample-specific and are not guaranteed to generalize to every benchmark
shift; current inferential uncertainty is not asserted until paired resampling
is regenerated from the retained prediction vectors.

% ---- Figure 4: expert complementarity (placeholder) -------------------------
\begin{figure}[htbp]
    \centering
    % TODO: \includegraphics[width=\linewidth]{figures/fig4_expert_complementarity.pdf}
    \caption{Complementarity and decision changes induced by expert fusion. (A) scatter of 1D vs SE3D probabilities; (B) error-overlap diagram; (C) probability shifts; (D) corrected/regressed/unchanged counts; (E, F) representative corrected cases.}
    \label{fig:expert_complementarity}
\end{figure}

\subsection{Analysis of the Task-Adaptive 1D Sequence Expert}
To investigate the role of the task-adaptive 1D expert, we first compared the
current sequence architecture with the earlier sequence pipeline based on
pooled ProtT5 representations and a TextCNN classifier
\cite{liang2026toxpltc}, and further examined
task adaptation, bidirectional state-space modeling, multi-seed ensembling,
and OOF-based probability calibration. Unlike approaches that use a pretrained
protein language model merely as a fixed feature extractor, GHXTox preserves
residue-level ProtT5 representations, adapts them to peptide toxicity
prediction through LoRA, and subsequently models residue dependencies in both
the N-to-C and C-to-N directions using bidirectional Mamba.

Importantly, the gain of the new 1D expert did not manifest as a uniform
increase across all threshold-dependent metrics. At their respective
OOF-frozen operating points, the previous ProtT5-TextCNN sequence candidate
achieved BACC, MCC, AUROC, and AUPRC values of 0.9317, 0.8586, 0.9724, and
0.9436, respectively, whereas the LoRA-ProtT5--Bi-Mamba expert achieved
0.9300, 0.8562, 0.9761, and 0.9533. Thus, BACC and MCC remained broadly
comparable, while AUROC and AUPRC increased by approximately 0.0037 and
0.0097, respectively.

This pattern indicates that the main advantage of the new sequence expert is
not simply a more aggressive decision boundary at one fixed threshold.
Instead, the task-adaptive representation improves the ordering of continuous
toxicity scores, particularly the precision--recall ranking of the positive
class. Such behavior provides a stronger probability stream for subsequent
fusion with structural evidence and highlights why sequence-model comparisons
should not rely on ACC or BACC alone.

\subsubsection{Task Adaptation of ProtT5 and the Retained 1D Representation}
Existing peptide-toxicity predictors commonly employ protein language models
to generate fixed sequence representations. For example, StrucToxNet uses
ProtT5-derived sequence features within its graph representation, whereas
PeptiTox integrates ESM2 embeddings with geometry derived from predicted
structures \cite{jiao2025structoxnet,wang2025peptitox}. A major distinction
of GHXTox is that ProtT5 is not used solely as
a static upstream encoder. Instead, its linear transformations are adapted to
the toxicity task through LoRA, while residue-level representations are
preserved for downstream sequence modeling.

This design serves two complementary purposes. First, low-rank updates allow
the pretrained semantic space to adapt to peptide toxicity prediction without
requiring full fine-tuning of the ProtT5 backbone. Second, peptide-level
pooling is deliberately postponed: contextualized residue states are passed
to the Bi-Mamba encoder, preserving position-dependent information and
allowing additional modeling of local and long-range residue dependencies.

Accordingly, the observed performance difference should not be attributed to
LoRA alone. We therefore interpret LoRA-ProtT5 and Bi-Mamba as a joint
task-adaptive sequence encoder. The current evidence supports the more
precise conclusion that task-adapted residue-level protein language model
representations followed by state-space sequence modeling improve probability
ranking while preserving threshold-dependent classification performance
comparable to that of the previous strong sequence baseline.

\subsubsection{Bidirectional State-Space Modeling of Peptide Sequences}
After task-adaptive ProtT5 encoding, GHXTox employs bidirectional Mamba
rather than TextCNN or immediate peptide-level pooling for secondary sequence
modeling. This choice is particularly relevant to peptides because
toxicity-related evidence may not be confined to a single local motif;
instead, combinations of charged, hydrophobic, aromatic, or other residues
may interact over distant sequence positions. Fixed-size convolutions can
capture such dependencies only through progressively enlarged receptive
fields, whereas a state-space model propagates contextual information through
recurrent state updates.

Although peptide sequences possess a natural N-to-C direction, toxicity
classification is not an autoregressive prediction task. GHXTox therefore
performs both forward and backward selective scans and integrates their
representations. Each residue can consequently be contextualized by
information from both termini rather than forcing the classifier to adopt a
single direction of information propagation.

A clear evidence boundary should nevertheless be maintained. The current
artifact set does not provide a strict single-variable comparison between
unidirectional and bidirectional Mamba under otherwise identical conditions.
We therefore do not assign a quantitative performance gain specifically to
bidirectionality. Instead, the evidence supports the complete
LoRA-ProtT5--Bi-Mamba configuration as a whole through group-aware OOF
evaluation and controlled candidate screening. This distinction separates an
architecture-level rationale from a causal ablation claim.

\subsubsection{Multi-Seed Ensembling, Calibration, and Controlled Negative Results}
After the basic sequence architecture had been established, we further
investigated model variance and probability reliability. The final 1D expert
was not obtained by selecting the random seed with the highest test
performance. Instead, three models were independently trained on the complete
training set and their probabilities were averaged with equal weights. A
larger ensemble constructed directly from the 15 fold-specific checkpoints
did not yield additional gains. A plausible explanation is that each fold
member is trained on only approximately 80\% of the available training
samples, and the resulting reduction in individual-model capacity outweighs
the variance reduction obtained from increasing the ensemble size. The final
configuration therefore uses the fixed average of three full-data models.

Following seed averaging, training OOF predictions were used to fit the
monotonic Platt calibrator. This operation primarily adjusts the probability
scale without altering the ranking of samples. The observed results were
consistent with this role: the 1D Brier score decreased from 0.0445 to 0.0429
on Test1 and from 0.0491 to 0.0451 on Test2. After subsequent fusion with the
3D expert, the Brier scores were further reduced to 0.0407 and 0.0430,
respectively. These results indicate that calibration is not merely a
cosmetic post-processing step, but an important component for producing
probability streams that can be meaningfully combined across experts.

Three intuitively attractive extensions were additionally evaluated but not
retained. Conservative positive-class sequence perturbation shifted
predictions toward the positive class and degraded specificity and the Brier
score. Length-specific calibration or thresholding was unstable because of
increased sampling variance within smaller strata and did not consistently
outperform the global operating point. Finally, the 15-member fold ensemble
did not outperform the three full-data models.

These negative results are methodologically informative. They indicate that,
at the current data scale, the limitation of the sequence expert cannot be
resolved simply by generating additional near-duplicate samples, increasing
the ensemble size, or fragmenting the decision rule into multiple
length-specific operating points. The retained three-seed full-data average
with OOF calibration represents a simpler and more stable solution and
reduces the risk of selective reporting or continued optimization against
historically observed test sets.

% ---- Table 5: 1D ablation ---------------------------------------------------
\begin{table}[htbp]
    \centering
    \caption{Selected controlled comparisons for the task-adaptive 1D sequence expert. OOF architecture rows and benchmark calibration rows are reported separately and should not be compared across evaluation splits.}
    \label{tab:1d_ablation}
    \resizebox{\linewidth}{!}{%
    \begin{tabular}{llccccc}
        \toprule
        Evaluation & Variant & BACC & MCC & AUROC & AUPRC & Brier \\
        \midrule
        Training OOF & Pooled ProtT5--TextCNN & 0.9317 & 0.8586 & 0.9724 & 0.9436 & NR \\
        Training OOF & LoRA-ProtT5--Bi-Mamba & 0.9300 & 0.8562 & \textbf{0.9761} & \textbf{0.9533} & NR \\
        \midrule
        Test1 & Three-seed 1D, uncalibrated & -- & -- & 0.9799 & 0.9644 & 0.0445 \\
        Test1 & Three-seed 1D, OOF Platt & -- & -- & 0.9799 & 0.9644 & \textbf{0.0429} \\
        Test2 & Three-seed 1D, uncalibrated & -- & -- & 0.9409 & 0.7831 & 0.0491 \\
        Test2 & Three-seed 1D, OOF Platt & -- & -- & 0.9409 & 0.7831 & \textbf{0.0451} \\
        \bottomrule
    \end{tabular}}
\end{table}

\subsection{Analysis of the pLDDT-Aware Multi-Scale 3D Structural Expert}
Compared with the 1D expert, the principal methodological innovations of
GHXTox lie in its structural expert. Previous studies have already
demonstrated the value of predicted structures for toxicity prediction.
PeptiTox combines ESM2 semantics with geometric descriptors extracted from
predicted structures, whereas StrucToxNet further employs an equivariant
graph neural network to preserve consistency under rigid transformations of
the three-dimensional coordinates
\cite{wang2025peptitox,jiao2025structoxnet}. For short peptides, however, a more
fundamental question is not simply whether three-dimensional coordinates are
available, but which structural evidence should be trusted and at which
spatial scale that evidence should be used.

The GHXTox structural expert is organized around this question. ESMFold
predictions are not treated as noise-free experimental structures. Instead,
pLDDT acts as a continuous reliability variable throughout geometric message
passing, residue-level gating, peptide-level gating, and training-time
coordinate perturbation. The model then extends beyond C$\alpha$-level
residue geometry by extracting local chemical pseudo-sites from full-atom PDB
structures and, when conformer ensembles are available, incorporating
statistics such as RMSF, distance variation, and contact occupancy through a
conservative residual pathway.

The central 3D contribution is therefore not merely the addition of more
structural features, but a hierarchical structural representation that
explicitly organizes evidence by reliability and spatial scale--from residue
geometry and local side-chain chemistry to sampled conformational variation.

\subsubsection{Continuous Structural-Confidence Modeling Is Preferable to Hard Filtering}
Structure-quality analysis first confirmed that pLDDT contains useful
information: peptides with higher mean pLDDT generally exhibit stronger
predictive performance. However, this does not imply that discarding
low-confidence samples is an optimal modeling strategy. In a controlled
high-pLDDT-only experiment, restricting training to high-confidence
structures substantially reduced performance over the complete test domain,
with MCC values of only 0.7232 on Test1 and 0.5593 on Test2.

This observation highlights an important distinction: structural confidence
is more useful for modulating structural evidence than for deciding whether a
sample should be retained at all. Low pLDDT may arise from genuine peptide
flexibility, limited information in short sequences, or uncertainty specific
to the structure-prediction model. Hard exclusion produces a structurally
cleaner subset but simultaneously removes many difficult examples that remain
part of the target prediction domain.

GHXTox therefore uses a continuous confidence-aware strategy. Low-confidence
edges remain in the graph but contribute weaker messages; residue-level gates
modulate structural evidence at individual positions; peptide-level gating
controls the overall use of structural context; and stronger coordinate
perturbations are applied to low-pLDDT residues during training.

This is an important distinction from approaches that analyze structure
quality only after model inference. In GHXTox, pLDDT is not merely a
diagnostic variable--it participates directly in representation learning.

\subsubsection{Full-Atom Chemical Sites Provide Consistent but Complementary Structural Gains}
A C$\alpha$-level residue graph effectively describes peptide compactness,
nonlocal contacts, and residue neighborhoods but compresses chemically
complex side chains into single residue nodes. To determine whether local
spatial chemistry provides additional discriminative information, we
introduced a full-atom-PDB-derived chemical-site residual on top of the
frozen C$\alpha$-EGNN. A strict paired five-fold design was adopted: in each
fold, the candidate model was initialized from the same control checkpoint,
while the sequence branch, C$\alpha$-EGNN, fusion module, and classifier were
frozen, leaving only the newly introduced chemical-site branch trainable.

Across the five CD-HIT homology-grouped folds, the complete chemical-site
model increased mean BACC from 0.9153 to 0.9171 ($\Delta=+0.0018$), F1 from
0.8790 to 0.8814 ($\Delta=+0.0023$), MCC from 0.8312 to 0.8342
($\Delta=+0.0030$), AUROC from 0.9687 to approximately 0.9692, and AUPRC from
0.9388 to 0.9401 ($\Delta=+0.0013$). More importantly, MCC, F1, and AUPRC
improved in all five folds.

The absolute magnitude of these gains is modest, but their consistency is
methodologically more informative. Because the base geometry, sequence
semantics, and classifier were frozen, the observed incremental performance
can only arise from the newly introduced chemical-site pathway under this
controlled training setup. The result therefore suggests that residue-centered
geometry already captures most of the structural signal, but local
side-chain chemistry contains a small yet reproducible component of
nonredundant predictive information.

Component-level ablations further support this interpretation. On the fixed
fold-0 screening subset, removing the second chemical site, orientation
invariants, either distance scale, explicit compatibility channels, pLDDT
weighting, or zero initialization decreased MCC. Retaining only the
raw-distance or only the normalized-distance representation reduced MCC by
approximately 0.0059 and 0.0058, respectively, whereas removing
interaction-type channels or pLDDT weighting each reduced MCC by
approximately 0.0057.

The methodological contribution is therefore not simply the construction of a
larger atom-level graph. GHXTox compresses local side-chain environments into
chemically interpretable pseudo-sites while simultaneously preserving
physical distances in Å, peptide-size-normalized distances, orientation
relationships, and candidate chemical compatibility.

PeptiTox demonstrated the utility of geometric information from predicted
structures, whereas ToxiPep highlighted complementarity between local atomic
chemistry and sequence context
\cite{wang2025peptitox,guan2025toxipep}. GHXTox integrates these ideas within a
confidence-aware 3D structural expert. The resulting interaction channels
should nevertheless be interpreted as candidate spatial-chemical
relationships available to the model, rather than confirmed hydrogen bonds,
salt bridges, cation--$\pi$ interactions, or disulfide bonds.

\subsubsection{Sequence--Structure Interaction and Hierarchical Confidence Gating}
Local geometric evidence becomes meaningful only in the appropriate sequence
context. The same spatial distance or hydrophobic contact may have different
implications depending on residue identity and sequence position. GHXTox
therefore avoids simply concatenating sequence and structural vectors after
peptide-level pooling. Instead, contextualized sequence states act as queries
that selectively retrieve evidence from pLDDT-gated spatial residue
representations.

This idea is conceptually related to the context--chemistry alignment
emphasized by ToxiPep, but in GHXTox the interaction occurs within a predicted
three-dimensional structural expert and is further regulated by both
residue-level and peptide-level structural-confidence gates.

The residue-level gate modulates how strongly the structural representation
of each residue contributes to the subsequent interaction, whereas the
peptide-level gate controls the overall contribution of structural context.
Structural information is therefore not switched off when confidence
decreases; rather, its contribution is continuously attenuated at multiple
levels.

Together with the negative result of hard high-pLDDT-only training described
in Section 3.4.1, these observations suggest that the value of the retained
design lies not merely in increasing model capacity, but in establishing a
learnable compromise between structural-information retention and
confidence-aware attenuation.

Because the current artifact set does not provide a complete single-variable
ablation of cross-attention and each confidence gate under otherwise identical
conditions, the overall structural gain is not quantitatively decomposed among
these components. The supported conclusion is therefore that continuous
pLDDT weighting, sequence--structure cross-attention, and hierarchical gating
jointly constitute the retained reliability-aware interaction mechanism.

\subsubsection{Conformer Statistics Provide Conditional Gains with Exact Fallback}
A single ESMFold prediction provides only one deterministic structural
anchor, whereas short peptides may exhibit substantial conformational
heterogeneity. GHXTox does not treat every generated conformer as an
additional labeled training example, nor does it repeatedly execute the
complete EGNN for every conformation. Instead, conformer ensembles are
compressed into statistics including RMSF, pairwise-distance variation,
contact occupancy at 6, 8, and 10~\AA, conformer RMSD, and radius of
gyration, and these statistics are allowed to produce only a zero-initialized
logit residual on top of the single-structure predictor.

The retained source-aware branch covers 2{,}795 of the 6{,}387 training
peptides. At the unchanged 0.72/0.28 fusion and decision threshold, activating
the full2795 residual increased training-OOF MCC from 0.87255 to 0.87445 and
BACC from 0.93220 to 0.93292, while reducing Brier score from 0.04316 to
0.04309. The MCC difference was $+0.00190$ under 5{,}000 class-stratified
paired bootstrap resamples, with a 95\% interval of
$[-0.00043,\,0.00454]$ and a paired-randomization $p$ value of 0.203.
AUROC and AUPRC decreased by 0.00014 and 0.00015, respectively. These results
support a small classification/calibration effect, but not a statistically
established or metric-uniform improvement.

The limited aggregate effect is informative. Conformer information is not
equally available across samples, and the two generators occupy different
length ranges. Moreover, source-specific scaling, confidence attenuation,
residual centering, generator-only activation, and residual clipping were all
evaluated with training-only nested OOF predictions and did not yield a
deployable improvement over the original source-aware residual. Additional
post hoc flexibility therefore did not convert broader conformer coverage
into a stable gain across all metric families.

The subsequent training-OOF information-ceiling analysis clarified this
failure mode. The sampled conformers were not geometrically redundant: the
median maximum pairwise RMSD was 4.786~\AA, and the median mean residue RMSF was
1.943~\AA. Nevertheless, the median range of the eight frozen-expert
probabilities was only $4.90\times10^{-6}$, with a 75th percentile of
$3.87\times10^{-4}$. The Pearson correlations between the probability range
and maximum pairwise RMSD and mean residue RMSF were 0.084 and 0.090,
respectively. None of the fixed label-free aggregators satisfied the frozen
promotion criteria. Even the nondeployable label-aware upper bound increased
full-cohort OOF MCC by only 0.00115 and corrected three decisions.

The internal spatial representations were more responsive than the final
probabilities: in the fold-0 extraction, the median pooled view-range norm was
1.350. However, within the later diagnostic's independently reproduced
single-structure baseline, the five-fold frozen spatial-set residual changed
full-cohort OOF MCC from 0.87077 to 0.87111
($\Delta=+0.00034$), while BACC, AUROC, AUPRC, and Brier score all moved in the
unfavorable direction. Its 95\% paired-bootstrap MCC interval was
$[-0.00084,\,0.00181]$, and Brier score worsened in 90.26\% of 5{,}000
bootstrap replicates. The candidate changed only three decisions, comprising
two corrections and one new error, and was therefore rejected without
benchmark evaluation or all-data refitting. These protocol-local values do not
replace or re-estimate the retained full2795 OOF comparison reported above.

Additional read-only analysis suggests that the principal bottleneck is not
the absence of distinguishable conformer representations, but their weak
alignment with the available toxicity labels. On the 2{,}795 covered records,
the 1D and base 3D probabilities were strongly correlated
($r=0.87084$). The learned residual had a correlation of $-0.0357$ with the
required fused-probability correction and moved the prediction toward the
observed label for only 45.87\% of covered records with a nonzero residual.
Thus the head changed its
output without learning a reproducible direction of correction, especially
for threshold-dependent metrics.

At frozen benchmark inference, the audited cache covers 486 Test1 peptides
and 284 Test2 peptides. The Test1 records comprise PepFlow conformers for
3--15 aa peptides and AlphaFold2/ColabFold conformers for 16--30 aa peptides;
Test2 coverage is obtained only through an exact unique-sequence match to the
audited Test1 cache. This routing rule was established independently of the
test labels.

For this reason, the conformer pathway is not implemented as a mandatory
backbone. Instead, GHXTox uses an availability-masked residual. For peptides
without valid conformer statistics, $\Delta z_{\mathrm{conf}} = 0$, and the
prediction therefore falls back exactly to the single-structure pathway.

The conformer module should therefore be interpreted as a conditional
structural refinement mechanism rather than a universal explanation for
performance gains across all benchmark samples.

\subsubsection{Structural Expert Performance and the Nature of the 3D Contribution}
After combining the retained components, the 3D-v2 structural expert achieved
BACC, MCC, AUROC, and AUPRC values of 0.9208, 0.8464, 0.9711, and 0.9474 on
Test1, compared with 0.9179, 0.8458, 0.9699, and 0.9426 for the historical
C$\alpha$-only 3D-v1 model. Test1 therefore showed a consistent but modest
improvement. Test2 exhibited a different trade-off: AUPRC increased from
0.6969 to 0.7241, whereas MCC decreased from 0.6320 to 0.6128.

This trade-off is particularly instructive. It indicates that finer-grained
structural modeling improves positive-class probability ranking but does not
guarantee simultaneous gains in all metrics at a frozen operating point. This
is particularly relevant for Test2, which contains only 46 toxic peptides,
where a small number of threshold-adjacent samples can substantially alter
MCC and F1.

The conclusion from the structural expert is therefore not that a more complex
3D model universally outperforms its predecessor. Rather, confidence-aware
multi-scale structural modeling introduces a small but consistent incremental
signal under group-aware training OOF evaluation and improves probability
ranking on selected historical benchmarks, while threshold-dependent gains
remain dataset-dependent.

This result also clarifies the role of the structural expert within GHXTox.
The 3D branch is most valuable as a source of complementary evidence rather
than as a universal replacement for a strong sequence predictor, providing a
direct rationale for the final dual-expert probability fusion.

% ---- Table 6: 3D ablation ---------------------------------------------------
\begin{table}[htbp]
    \centering
    \caption{Ablation analysis of the 3D structural expert. The first four rows use historical benchmark predictions; the final two rows are a paired five-fold OOF chemical-site experiment.}
    \label{tab:3d_ablation_complexity}
    \resizebox{\linewidth}{!}{%
    \begin{tabular}{llcccc}
        \toprule
        Evaluation & Variant & BACC & MCC & AUROC & AUPRC \\
        \midrule
        Test1 & Historical C$\alpha$ 3D-v1 & 0.9179 & 0.8458 & 0.9699 & 0.9426 \\
        Test1 & pLDDT-aware multi-scale 3D-v2 & \textbf{0.9208} & \textbf{0.8464} & \textbf{0.9711} & \textbf{0.9474} \\
        Test2 & Historical C$\alpha$ 3D-v1 & \textbf{0.8481} & \textbf{0.6320} & \textbf{0.9316} & 0.6969 \\
        Test2 & pLDDT-aware multi-scale 3D-v2 & 0.8453 & 0.6128 & 0.9295 & \textbf{0.7241} \\
        \midrule
        Training OOF & Frozen backbone, no chemical-site residual & 0.9153 & 0.8312 & 0.9687 & 0.9388 \\
        Training OOF & Frozen backbone + chemical-site residual & \textbf{0.9171} & \textbf{0.8342} & \textbf{0.9692} & \textbf{0.9401} \\
        \bottomrule
    \end{tabular}}
\end{table}

% ---- Figure 5: module-level ablations (placeholder) --------------------------
\begin{figure}[htbp]
    \centering
    % TODO: \includegraphics[width=\linewidth]{figures/fig5_module_ablation.pdf}
    \caption{Module-level ablation results for the 1D and 3D experts (A--F).}
    \label{fig:module_ablation}
\end{figure}

\subsection{Robustness, Statistical Uncertainty, and Applicability Analysis}
Point estimates alone are insufficient for distinguishing a stable model
improvement from variation caused by finite samples. We therefore retain
multi-seed averaging for the 1D expert and use paired, sample-level bootstrap
resampling to compare the source-aware residual with its same-checkpoint
zero-conformer OOF baseline. All comparisons preserve the pairing between the
two predictions for each training peptide.

\subsubsection{Multi-Seed Reproducibility}
Three full-training 1D models initialized with seeds 42, 123, and 2025 are
averaged before calibration. This ensemble is the deployed predictor; no seed
is selected according to Test1 or Test2 performance. A 15-member ensemble of
fold-specific checkpoints was also examined but not retained because its
members had each seen only approximately 80\% of the training set and the
aggregate did not improve the frozen validation criteria.

\subsubsection{Paired Uncertainty and Point-Estimate Comparison}
The current training-OOF comparison uses 5{,}000 class-stratified paired
bootstrap resamples and 5{,}000 within-sample randomizations. Relative to the
same-checkpoint zero-conformer OOF vector, the MCC difference is
$+0.00190$ (95\% CI $[-0.00043,\,0.00454]$; randomization
$p=0.203$), and the BACC difference is $+0.00071$ (95\% CI
$[-0.00038,\,0.00181]$). Thus the point estimates favor the source-aware
residual at the frozen threshold, but the intervals include zero.

The archived per-sample OOF vector for the immediately preceding pilot
checkpoint is not available. Table~\ref{tab:bootstrap_paired} therefore
reports deterministic benchmark differences between the source-aware
full2795 model and that preceding pilot default, without attaching a
full2795-versus-pilot confidence interval. These retrospective differences
are descriptive and were not used to tune the source-aware checkpoint.

% ---- Table 7: paired point-estimate statistics ------------------------------
\begin{table}[htbp]
    \centering
    \caption{Point-estimate changes from the preceding pilot default to the retained source-aware full2795 model. Direct paired confidence intervals are not reported because the pilot per-sample OOF vector is unavailable.}
    \label{tab:bootstrap_paired}
    \begin{tabular}{llcc}
        \toprule
        Dataset & Metric & Current--previous $\Delta$ & Favorable direction \\
        \midrule
        \multirow{4}{*}{Test1}
        & BACC & $-0.00094$ & Increase \\
        & MCC & $-0.00018$ & Increase \\
        & AUPRC & +0.00006 & Increase \\
        & Brier & $-0.00030$ & Decrease \\
        \midrule
        \multirow{4}{*}{Test2}
        & BACC & +0.00093 & Increase \\
        & MCC & +0.00725 & Increase \\
        & AUPRC & $-0.00065$ & Increase \\
        & Brier & $-0.00082$ & Decrease \\
        \bottomrule
    \end{tabular}
\end{table}

\subsubsection{Controlled Multi-Conformer Ablation on Test1 and Test2}
The contribution of the conformer pathway was isolated with a same-checkpoint
ablation. The LoRA-ProtT5--Bi-Mamba probabilities, structural checkpoint,
0.72/0.28 fusion weights, and threshold of 0.47022256531284556 were held fixed;
only the audited conformer statistics were either supplied or replaced by the
exact availability-masked zero-conformer input. This design attributes the
reported differences specifically to activation of the conformer residual,
rather than to retraining, recalibration, or operating-point selection.

% ---- Reserved Test1 multi-conformer ablation table position -----------------
\begin{table}[p]
    \centering
    \caption{Controlled multi-conformer ablation on Test1. The zero-conformer and multi-conformer variants share the same trained checkpoint, 1D probabilities, fusion weights, and decision threshold. The multi-conformer cache covers 486 of 1{,}126 Test1 peptides. $\Delta$ is computed as multi-conformer minus zero-conformer; positive values are favorable for all listed metrics except Brier score, for which a negative value is favorable.}
    \label{tab:conformer_ablation_test1}
    \begin{tabular}{lrrr}
        \toprule
        Metric & Zero conformer & Multi-conformer & $\Delta$ \\
        \midrule
        ACC       & 0.947602 & 0.948490 & +0.000888 \\
        BACC      & 0.932308 & 0.932929 & +0.000620 \\
        Precision & 0.916933 & 0.919872 & +0.002939 \\
        SN        & 0.896875 & 0.896875 & +0.000000 \\
        SP        & 0.967742 & 0.968983 & +0.001241 \\
        F1        & 0.906793 & 0.908228 & +0.001435 \\
        MCC       & 0.870460 & 0.872568 & +0.002108 \\
        AUROC     & 0.981056 & 0.981099 & +0.000043 \\
        AUPRC     & 0.967403 & 0.967482 & +0.000080 \\
        Brier     & 0.040991 & 0.040745 & $-0.000246$ \\
        \bottomrule
    \end{tabular}
\end{table}

On Test1, activation of the audited source-aware residual changes one
thresholded prediction, converting one false positive into a true negative.
Consequently, SP, precision, BACC, F1, and MCC increase, whereas SN is
unchanged. AUROC and AUPRC also increase slightly, and Brier score decreases
by approximately $2.46\times10^{-4}$. The gains are directionally consistent
but small in absolute magnitude.

% ---- Reserved Test2 multi-conformer ablation table position -----------------
\begin{table}[p]
    \centering
    \caption{Controlled multi-conformer ablation on Test2. The zero-conformer and multi-conformer variants share the same trained checkpoint, 1D probabilities, fusion weights, and decision threshold. The multi-conformer cache covers 284 of 582 Test2 peptides through an audited exact unique-sequence join to Test1. $\Delta$ is computed as multi-conformer minus zero-conformer; positive values are favorable for all listed metrics except Brier score, for which a negative value is favorable.}
    \label{tab:conformer_ablation_test2}
    \begin{tabular}{lrrr}
        \toprule
        Metric & Zero conformer & Multi-conformer & $\Delta$ \\
        \midrule
        ACC       & 0.946735 & 0.948454 & +0.001718 \\
        BACC      & 0.861778 & 0.862711 & +0.000933 \\
        Precision & 0.636364 & 0.648148 & +0.011785 \\
        SN        & 0.760870 & 0.760870 & +0.000000 \\
        SP        & 0.962687 & 0.964552 & +0.001866 \\
        F1        & 0.693069 & 0.700000 & +0.006931 \\
        MCC       & 0.667340 & 0.674587 & +0.007248 \\
        AUROC     & 0.943300 & 0.943300 & +0.000000 \\
        AUPRC     & 0.804590 & 0.804328 & $-0.000262$ \\
        Brier     & 0.043631 & 0.042987 & $-0.000644$ \\
        \bottomrule
    \end{tabular}
\end{table}

On Test2, the source-aware residual also converts one false positive into a
true negative. MCC increases by 0.00725, F1 by 0.00693, and Brier score
decreases by approximately $6.44\times10^{-4}$, while SN is unchanged.
AUROC is unchanged and AUPRC decreases by $2.62\times10^{-4}$. Across the two
benchmarks, conformer statistics therefore make a conditional and
quantitatively modest contribution, with a classification/calibration benefit
but no uniform improvement across metric families.

\subsection{Interpretability and Representative Case Studies}
All interpretation experiments used the retained frozen model and the
prespecified case-selection rule described above. The analysis included 460
residue occlusions across 16 peptides, unperturbed confidence gates for all
1{,}708 benchmark records, 445 individual chemical-site deletions, 121
category-level deletions, and 424 residue-level site deletions. Before
perturbation, baseline probabilities were recomputed and compared with the
archived prediction vectors. The maximum absolute discrepancy for the 1D
expert within the case cohort was $3.45\times10^{-4}$; the corresponding 3D
discrepancies were $5.49\times10^{-4}$ on Test1 and $1.14\times10^{-5}$ on
Test2. None of these differences changed a standalone 3D or fused decision.
The perturbation baselines therefore preserved the archived decision states,
with the larger Test1 discrepancy reported as bounded cross-environment
numerical drift.

\subsubsection{Residue-Level Interpretation of the 1D Expert}
The residue-occlusion profiles were heterogeneous across peptides and seeds.
Across the 16 prespecified cases, the mean pairwise Spearman correlation of
absolute per-residue logit effects was 0.118 (case-level range, $-0.077$ to
0.422), whereas the mean residue-wise sign-agreement fraction was 0.740
(range, 0.667--0.796). Direction was therefore more stable across seeds than
the fine-grained ranking of residue magnitudes. The reserved maps emphasize
local effects with consistent signs rather than implying a single stable
ranking across all residues.

Several cases nevertheless contained large, directionally consistent effects.
For the Test2 toxic peptide \texttt{sequence\_34\_34}, occlusion of Asn9
reduced the calibrated 1D probability from 0.5447 to 0.0412
($\Delta p_{1D}=+0.5035$), with all three seeds agreeing on the sign. In the
Test1 toxic case \texttt{peptide\_199}, occlusion of Ser24 reduced the score
from 0.3687 to 0.1114 ($\Delta p_{1D}=+0.2573$), again with complete sign
agreement. Effects were not uniformly toxicity-supporting: in the Test2 toxic
case \texttt{sequence\_23\_23}, occlusion of Ile5 increased the 1D probability
from 0.2807 to 0.4819 ($\Delta p_{1D}=-0.2013$). These cases illustrate large,
context-dependent sensitivities in the 1D output; they do not identify
experimentally validated toxic residues. The dataset-specific layouts are
reserved in
Figures~\ref{fig:test1_residue_occlusion} and
\ref{fig:test2_residue_occlusion}, respectively.

% ---- Reserved figure: Test1 residue occlusion -------------------------------
\begin{figure}[H]
    \centering
    \fbox{\parbox[c][0.22\textheight][c]{0.96\linewidth}{
        \centering
        \textbf{Reserved for the Test1 residue-occlusion figure}\\[6pt]
        \small Eight sequence heat strips will show one toxic and one non-toxic
        peptide from each expert-outcome group. A shared signed color scale will
        encode within-peptide $L_1$-normalized $\Delta p_{1D}$; residue positions,
        seed sign agreement, baseline probability, and expert-outcome group will
        be annotated. Cases are ordered by the prespecified cohort rule, not by
        attribution magnitude.
    }}
    % TODO: \includegraphics[width=\linewidth]{figures/test1_residue_occlusion.pdf}
    \caption{Reserved layout for Test1 residue-level analysis of the frozen 1D expert.
    Each row will represent one independently selected peptide and each cell one
    residue. Positive values will indicate that zeroing the post-contextual
    residue vector lowers the calibrated 1D toxicity probability, whereas
    negative values will indicate an increase. The common color normalization
    will permit comparison of direction and relative within-peptide magnitude;
    it will not imply that equal colors correspond to equal biological effects
    across peptides. Seed agreement will distinguish reproducible directions
    from effects dominated by one checkpoint.}
    \label{fig:test1_residue_occlusion}
\end{figure}

% ---- Reserved figure: Test2 residue occlusion -------------------------------
\begin{figure}[H]
    \centering
    \fbox{\parbox[c][0.22\textheight][c]{0.96\linewidth}{
        \centering
        \textbf{Reserved for the Test2 residue-occlusion figure}\\[6pt]
        \small The layout, case ordering, labels, and signed normalization will
        match the Test1 panel exactly. Separate presentation preserves benchmark
        provenance while retaining a directly comparable visual design.
    }}
    % TODO: \includegraphics[width=\linewidth]{figures/test2_residue_occlusion.pdf}
    \caption{Reserved layout for Test2 residue-level analysis of the frozen 1D expert.
    The eight rows will cover both labels in all four expert-outcome groups.
    Residue-wise probability shifts will be computed with the three frozen seeds,
    the original deployment padding context, and the fixed OOF Platt transform.
    Test2 comprises exact sequence matches to a subset of Test1, so matched
    records may yield similar maps. Such agreement will be treated as a
    reproducibility check, not as independent external validation.}
    \label{fig:test2_residue_occlusion}
\end{figure}

\subsubsection{Structural Confidence and Chemical-Site Sensitivity}
Mean pLDDT was strongly associated with the learned global structural gate.
Spearman correlations were
0.960 on Test1 (95\% peptide-bootstrap interval, 0.953--0.966) and 0.950 on
Test2 (0.937--0.959); the pooled descriptive value was 0.957
(0.951--0.961). In contrast, the mean within-peptide correlation between
residue pLDDT and node gates was small: 0.049 on Test1 (0.030--0.068) and
0.064 on Test2 (0.036--0.092). Thus, global gating closely tracked confidence
at the peptide level, whereas the weak local correlations are consistent with
residue gates also reflecting geometric and contextual features.

Most individual-site deletions produced small changes, but the distribution
had a pronounced upper tail. Across 445 deletions, the median and mean absolute changes in
$p_{3D}$ were 0.00048 and 0.00997, respectively; the 95th percentile was
0.0570 and the maximum was 0.2994. Nine individual deletions moved the
recomputed fused probability across the frozen decision threshold when the 1D
probability was held fixed. At category level, hydrophobic-site deletion
had the largest mean absolute change among the eight modeled categories
(0.126 across 16 case-level deletions), followed by aromatic sites (0.0687),
whereas sulfur-site deletion had the smallest mean absolute change (0.00179).
Because the analysis used the selected 16-peptide cohort, these category-level
estimates are descriptive rather than population effects.

The largest individual-site effect occurred at the aromatic/hydrophobic Phe8
site in the matched cache-covered cases \texttt{peptide\_846} and
\texttt{sequence\_388\_388}: deletion increased $p_{3D}$ from approximately
0.2405 to 0.5399, giving $\Delta p_{3D}\approx-0.2994$. Conversely, deletion
of the hydrophobic/conformational Pro19 site in \texttt{peptide\_199}
decreased $p_{3D}$ from 0.3443 to 0.1163
($\Delta p_{3D}=+0.2280$). The opposing directions show that the deletion
score depends on the surrounding peptide context and cannot be read as an
intrinsic toxicity label for a functional group. Figure~\ref{fig:structural_sensitivity}
reserves a layout for these scale-specific associations and perturbation effects.

% ---- Reserved figure: structural interpretation ----------------------------
\begin{figure}[H]
    \centering
    \fbox{\parbox[c][0.28\textheight][c]{0.96\linewidth}{
        \centering
        \textbf{Reserved for the structural-sensitivity figure}\\[6pt]
        \small (A) Mean pLDDT versus the learned global gate for Test1 and Test2,
        with dataset-specific Spearman coefficients and bootstrap intervals.
        (B) Distribution of within-peptide pLDDT--node-gate correlations.
        (C) Signed and absolute $\Delta p_{3D}$ for category-level deletions.
        (D) Largest individual-site deletions, labeled by residue, position,
        chemical-site type, conformer availability, and induced fused shift.
    }}
    % TODO: \includegraphics[width=\linewidth]{figures/structural_sensitivity.pdf}
    \caption{Reserved layout for structural-confidence and chemical-site
    sensitivity analysis. (A) Association between mean pLDDT and the global
    gate. (B) Distribution of within-peptide pLDDT--node-gate correlations.
    (C) Comparison of the eight modeled chemical categories. (D) Largest
    probability-increasing and probability-decreasing site-deletion effects.
    The deletion panels use the frozen 3D expert and calculate
    $\Delta p_{\mathrm{fused}}$ with $p_{1D}$ fixed; the reported values denote
    model sensitivity, not experimentally verified toxicophores.}
    \label{fig:structural_sensitivity}
\end{figure}

\subsubsection{Expert Agreement, Disagreement, and Fusion Outcomes}
The two experts were both correct for 1{,}032 of 1{,}126 Test1
peptides (91.65\%) and 529 of 582 Test2 peptides (90.89\%). On Test1, the 1D
expert alone was correct for 29 cases (2.58\%), the 3D expert alone for 23
(2.04\%), and both were incorrect for 42 (3.73\%). The corresponding Test2
counts were 13 (2.23\%), 14 (2.41\%), and 26 (4.47\%). The median absolute
expert-probability gaps were small in the both-correct groups (0.0145 on Test1
and 0.0126 on Test2) but exceeded 0.49 in each one-expert-correct group. The
latter cases therefore reflect substantial disagreement, not threshold-scale
numerical variation.

Among samples for which only one expert was correct, the fixed 0.72/0.28
fusion remained correct in 19 of 29 1D-only-correct Test1 cases and 17 of 23
3D-only-correct cases. The corresponding Test2 counts were 10 of 13 and 12 of
14. For example, the toxic Test2 peptide \texttt{sequence\_23\_23} had
$p_{1D}=0.2806$, $p_{3D}=0.9954$, and $p_{\mathrm{fused}}=0.4808$; the high
3D score moved the fused prediction above the decision threshold. By contrast,
the toxic
\texttt{sequence\_34\_34} had $p_{1D}=0.5447$ but
$p_{3D}=3.0\times10^{-5}$, yielding $p_{\mathrm{fused}}=0.3922$ and a false
negative despite the correct standalone 1D decision. These contrasting cases
show that structural evidence can either correct or introduce an error under a
fixed weighted fusion. Figure~\ref{fig:expert_agreement_cases} reserves a common
layout for the outcome counts, probability disagreements, and cache-covered
case studies.

% ---- Reserved figure: expert agreement and cases ---------------------------
\begin{figure}[H]
    \centering
    \fbox{\parbox[c][0.27\textheight][c]{0.96\linewidth}{
        \centering
        \textbf{Reserved for the expert-agreement and case-study figure}\\[6pt]
        \small (A) Outcome-group proportions for Test1 and Test2.
        (B) $p_{1D}$ versus $p_{3D}$ with frozen expert thresholds and the
        fusion decision boundary. (C) Unified probability bars for the 16
        independently selected cases. (D) Matched cache-covered records across
        Test1 and Test2, with representative site-deletion shifts and explicit
        availability masks.
    }}
    % TODO: \includegraphics[width=\linewidth]{figures/expert_agreement_cases.pdf}
    \caption{Reserved layout for expert agreement, disagreement, and
    representative cases. (A) Proportions of all four predeclared outcome
    groups. (B) Magnitude and direction of expert disagreement. (C) Probability
    profiles for the label-balanced, threshold-proximal perturbation cohort,
    annotated by whether fusion agrees with the individually correct expert.
    (D) Matched, cache-covered cases shown as sensitivity analyses of
    predicted-structure inputs, not as experimental conformers or
    molecular-dynamics trajectories.}
    \label{fig:expert_agreement_cases}
\end{figure}

\subsubsection{Representative Cases with Audited Conformer Coverage}
The interpretability cohort contains cases with audited conformer-cache
coverage and cases assigned the prespecified zero-residual fallback. Among the
445 individual-site perturbations, 144
belonged to cache-covered peptides and 301 to fallback peptides. Their mean
absolute $\Delta p_{3D}$ values were 0.0166 and 0.00678, respectively. This
unmatched comparison cannot isolate an effect of conformer availability because
the cases differ in sequence, label, length, site count, and baseline confidence.

Several selected Test2 peptides exactly match sequences selected from Test1 and
therefore provide a cross-dataset reproducibility check for preprocessing and
perturbation.
For the cache-covered pair \texttt{peptide\_481}/
\texttt{sequence\_143\_143}, deleting the donor/acceptor Asn18 site changed
$p_{3D}$ from approximately 0.4653 to 0.6420 in both benchmark records
($\Delta p_{3D}\approx-0.1767$), and deleting the hydrophobic Ile6 site gave
$\Delta p_{3D}\approx-0.1431$ in both. The near-identical values are consistent
with the declared sequence-based reuse of the Test2 cache and the deterministic
perturbation procedure. They are not independent replicates because Test2
comprises exact sequence matches to a subset of Test1. No conclusions about
conformational dynamics are drawn because aligned trajectories, RMSF, contact
occupancy, and experimentally determined conformational transitions were not
analyzed.

\subsection{Limitations and Future Directions}
GHXTox combines task-adaptive sequence modeling, reliability-aware
three-dimensional structural representation, and OOF-frozen probability
fusion. The present evidence is nevertheless bounded by data provenance,
predicted-structure uncertainty, conformer coverage, the fixed fusion strategy,
and the scope of the post hoc analyses.

Structural extensions, dynamic fusion schemes, and conformer-expansion
candidates that did not meet the frozen promotion criteria were retained in the
audit trail. Their mixed results suggest that additional architectural
complexity alone will not resolve the current data and validation constraints.

\subsubsection{Dataset and Evaluation Limitations}
Although calibration, fusion parameters, and the decision threshold were
determined using group-aware OOF predictions without any formal tuning on
Test1 or Test2, both benchmark sets had been observed during historical model
development. They should therefore be regarded as frozen historical
benchmarks rather than newly blinded external test sets. In addition, Test2
contains only 46 toxic peptides, making SN, BACC, and MCC sensitive to a
small number of prediction changes. Stronger evidence of cross-domain
generalization will require a new external dataset evaluated under a prefrozen
protocol.

\subsubsection{Structural and Conformer Limitations}
The 3D expert depends on ESMFold-predicted structures and is therefore
subject to structure-prediction uncertainty. Although pLDDT allows structural
evidence to be continuously down-weighted, it is not equivalent to
experimental structural accuracy, conformational flexibility, or dynamical
stability. Similarly, full-atom-derived chemical sites represent candidate
spatial-chemical relationships available to the model rather than confirmed
hydrogen bonds, salt bridges, or cation--$\pi$ interactions.

Conformer coverage is also uneven: audited ensembles are available for
43.16\% of Test1 and 48.80\% of Test2, and generator identity is partially
coupled to peptide length because PepFlow supplies the retained 3--15 aa
records whereas AlphaFold2/ColabFold supplies the retained 16--30 aa records.
The conformer pathway is therefore a conditional refinement when ensemble
information is available, not a universal source of improvement. Future work
should construct same-sample,
multi-generator conformer sets to separate generator effects from
peptide-length effects and should evaluate their calibration consequences on a
new prefrozen external dataset.

The present residual also remains deliberately simple. Training-only nested
OOF experiments with source-specific scaling, confidence attenuation,
source-wise centering, generator-only activation, and logit clipping did not
improve the deployed full2795 predictions. More direct aggregation did not
resolve the problem: although the available conformers were geometrically
diverse and their internal spatial embeddings varied, the frozen structural
expert produced nearly invariant view-level probabilities, and a compact
spatial-set residual did not provide a stable multimetric OOF benefit. The
retained full2795 pathway should therefore be regarded as a qualified
cache-enabled extension with exact single-structure fallback, whereas the later
frozen spatial-set candidate is nondeployable.

These findings identify a representation--label alignment limitation rather
than demonstrating that peptide conformational heterogeneity is biologically
irrelevant. Each peptide has a single sample-level toxicity label, without an
experimentally identified active conformer, conformer population, assay
condition, or mechanism-specific endpoint. Consequently, the current objective
does not tell the model which structural transition or microstate should carry
the label. Direct paired uncertainty against the preceding pilot default also
cannot be computed because its per-sample OOF vector and fold checkpoints are
unavailable.

\subsubsection{Modeling and Interpretability Limitations}
GHXTox currently trains the two experts separately and combines their
probabilities by linear fusion rather than full end-to-end joint
optimization. This design provides clear training and calibration boundaries,
but a fixed fusion weight cannot adapt the relative importance of sequence
and structural evidence for individual peptides. Dynamic gating and stacking
strategies were explored but did not consistently outperform the frozen
configuration, suggesting that greater fusion flexibility may introduce
additional overfitting at the current data scale.

Interpretability remains incomplete despite the matched analyses reported
here. The 1D experiment zeros post-contextual residue vectors in a selected
16-peptide cohort and therefore does not measure the effect of a realizable
amino-acid substitution or capture changes that would occur during ProtT5
contextualization. Its modest across-seed rank agreement further limits claims
about a unique residue ordering. The chemical-site deletions operate on
predicted-structure features, and the strong mean-pLDDT--global-gate
association partly reflects the fact that confidence is an explicit model
input. Neither analysis validates a causal toxicity mechanism, an experimental
conformer, or a biochemical intervention.

\subsubsection{Future Directions}
Evaluation on a new independent external dataset is the immediate priority for
establishing stronger generalization evidence. Expanding conformer coverage
will also require explicit control of peptide length and structure generator.
Future structural representations may incorporate solvent environment,
membrane context, protonation state, and peptide topology.

For multi-conformer modeling, the next step should target representation
learning rather than add further fixed aggregators to frozen probabilities. A
bounded, fold-pure experiment could partially unfreeze only the terminal
spatial and chemical layers while preserving the single-structure anchor, and
could add auxiliary objectives for residue RMSF, contact-occupancy changes, or
inter-conformer geometric consistency. Residue-wise microstate fusion may also
retain local transitions that disappear under peptide-level pooling. Such a
candidate should proceed to benchmark evaluation only after demonstrating a
consistent OOF benefit across MCC, BACC, AUPRC, and Brier score, without
degrading the exact fallback pathway.

If sufficiently reliable metadata become available, mechanism-aware endpoints
and assay context should be modeled explicitly. Labels that distinguish, for
example, membrane disruption, hemolysis, and general cytotoxicity may provide a
more informative supervisory signal for identifying relevant conformational
states than a single pooled toxicity label. Experimentally supported conformer
populations or environment-specific ensemble weights would further reduce the
current ambiguity between sampled structural diversity and biologically active
microstates.

The fixed-representation occlusion analysis should be compared with orthogonal
token-level perturbations and gradient-based attribution on a larger
prespecified cohort. Hypotheses supported by both the 1D and 3D analyses could
then be evaluated, where feasible, through controlled mutational, binding,
membrane-disruption, or toxicity experiments.

\section{CONCLUSIONS}
GHXTox combines a task-adaptive LoRA-ProtT5--Bi-Mamba sequence expert with a
pLDDT-aware multi-scale structural expert through an auditable probability
fusion protocol. The structural expert integrates residue geometry,
full-atom-derived chemical sites, and a source-aware, availability-masked
conformer residual, while the final predictor uses the structural expert's
standalone probability rather than recycling a previous fused output. The
1D/3D weights (0.72/0.28) and the
decision threshold (0.47022256531284556) were fixed from training OOF
predictions before benchmark evaluation.

On the frozen historical Test1 and Test2 benchmarks, the current model reaches
MCC values of 0.8726 and 0.6746 and AUPRC values of 0.9675 and 0.8043,
respectively. Relative to the calibrated 1D probability stream, the final
prespecified fusion yielded higher AUROC and AUPRC values and lower Brier scores
on both sets, consistent with complementary structural information. The
controlled conformer ablation remains modest: it corrects one false positive
on each benchmark and reduces Brier score on both, but Test2 AUPRC decreases
slightly. In training OOF, the MCC difference from the same-checkpoint
zero-conformer predictor is $+0.00190$, with a 95\% paired-bootstrap interval
of $[-0.00043,\,0.00454]$. The point estimates therefore do not support claims
of statistically established or universal superiority. A direct paired
full2795-versus-pilot OOF interval also remains unavailable because the
preceding pilot prediction vector was not retained. The training-OOF conformer
diagnostics further show that marked geometric and embedding variation can
coexist with nearly invariant frozen-expert probabilities. The rejected
spatial-set residual did not produce a stable multimetric gain, indicating a
representation--label alignment bottleneck under the current sample-level
supervision rather than establishing that conformational heterogeneity lacks
biological relevance. The matched interpretation experiments identified a
strong peptide-level association
between predicted structural confidence and the global gate, mostly small but
occasionally large chemical-site deletion effects, and modest cross-seed
stability in the fine-grained 1D residue ranking. These results describe model
sensitivity rather than causal mechanisms. Future validation should prioritize
a genuinely new, prefrozen external dataset, better-controlled same-sample
conformer evidence, comparisons with orthogonal attribution methods, and
experimental evaluation of the resulting hypotheses.

% ----------------------------------------------------------------------------
%  Unnumbered sections after the conclusions (as in the reference paper)
% ----------------------------------------------------------------------------
\section*{Data and Software Availability}
The source code, configuration files, fold manifests, and inference scripts
used in this study are available at
\url{https://github.com/Suiyumanbu/GHXTox}. A versioned archival release
containing the exact commit identifier, environment manifest, frozen model
parameters, calibration coefficients, fusion weights, full-precision decision
threshold, and sample-level prediction files will accompany the final
publication.
% TODO: freeze the final commit, release/tag, archive DOI, license, and
%       environment files before submission.
Raw peptide records were obtained from the public databases described in
Section 2.1. Where redistribution is restricted by the original database
terms, accession lists and preprocessing scripts are provided instead of
redistributing the source records.
% TODO: processed data and fold manifests (train/test manifest, source_index,
%       sequence, label, homology group, fold); three 1D checkpoints and one
%       final 3D checkpoint; Platt scale/bias; fusion weights; full-precision
%       threshold; Test1/Test2 prediction CSVs; metric JSON; bootstrap and
%       statistical outputs. Note reproduction limits or third-party licenses.

\section*{Supporting Information}
A Supporting Information package containing detailed dataset audits, fold-wise
results, hyperparameter settings, additional ablation experiments, subgroup
analyses, statistical outputs, and extended interpretability case studies will
accompany the submitted manuscript.
% TODO: provide the supplementary PDF/data pack (e.g., Supplementary Methods
%       S1--S6, Results S7--S16, Tables S1--S12, Figures S1--S10).

\section*{CRediT Authorship Contribution Statement}
% TODO: assign CRediT roles per author, e.g.:
% Author A: Conceptualization, Methodology, Software, Formal analysis,
%           Writing -- original draft.
% Author B: Data curation, Validation, Visualization, Software.
% Author C: Structural modeling, Investigation, Interpretation,
%           Writing -- review and editing.
% Author D: Supervision, Project administration, Funding acquisition,
%           Writing -- review and editing.

\section*{Acknowledgements}
% TODO: acknowledge supervisors, computing resources or platforms, and
%       contributors who do not meet authorship criteria.

\section*{Declaration of Competing Interest}
The authors declare that they have no known competing financial interests or
personal relationships that could have appeared to influence the work
reported in this paper.
% TODO: update if any patent, company, consulting, or equity relationship exists.

% ============================================================================
%  References: numbered by first appearance in the text.
% ============================================================================
\bibliographystyle{unsrt}
\bibliography{references}

\end{document}
